\documentclass[
    12pt
    ,oneside
    ,a4paper
    ,chapter=TITLE
    ,section=TITLE
    ,sumario=abnt-6027-2012]{abntex2}
    
    % Carrega o pacote da abnt com as adapatções para normas da Tuiuti
\usepackage{../packages/abnt-UTP}
\usepackage[alf]{abntex2cite}

\renewcommand{\tt}{\ttfamily}

% =================================================================================================
% -------------------------------------------------------------------------------------------------
% -                                                                                               -
% - INFORMAÇÕES GERAIS                                                                            -
% -                                                                                               -
% -------------------------------------------------------------------------------------------------
% =================================================================================================

% Informações para CAPA e FOLHA DE ROSTO
\titulo{SISTEMA INTEGRADO DE MONITORAMENTO EM SERVIDOR LINUX:\\
        MONITORAMENTO DE INFRAESTRUTURA E EVENTOS DE SEGURANÇA}

\autor{Beatriz Pimentel de Mello \\
       Caio Federico Esquivel Lovera Arze\\  
       Denyse Panza Clemente Ferreira\\
       Douglas Clayton da Silva \\
       Fernando Bach \\
       Larissa Quirino dos Santos \\
       Lucas Toterol Rodrigues
       }

\orientador{Prof. Luiz Altamir Corrêa Junior}

\preambulo{Projeto apresentado ao curso de Análise e Desenvolvimento de Sistemas, da Universidade Tuiuti do Paraná, como requisito avaliativo da disciplina de Projeto Interdisciplinar: Análise de Sistemas}

\instituicao{Universidade Tuiuti do Paraná}
\local{Curitiba}
\data{\the\year}


\begin{document}

% =================================================================================================
% -------------------------------------------------------------------------------------------------
% - ELEMENTOS PRÉ-TEXTUAIS                                                                        -
% -------------------------------------------------------------------------------------------------
% =================================================================================================

% -------------------------------------------------------------------------------------------------
% Cria a capa e a folha de rosto
\imprimircapa
\imprimirfolhaderosto

\begin{resumo}
    A crescente dependência de ambientes computacionais torna a indisponibilidade e os incidentes de segurança fatores críticos para a continuidade dos serviços. Este trabalho apresenta a fundamentação teórica, o planejamento arquitetural e a implementação de um sistema de monitoramento para servidores Linux, com foco na coleta de métricas de infraestrutura e na detecção de eventos básicos de segurança por meio da análise de logs. Com base em conceitos de monitoramento, observabilidade e segurança da informação, foi implementada uma arquitetura cliente-servidor na qual agentes coletores instalados nos dispositivos monitorados capturam dados de CPU, memória, disco e rede e os enviam a um servidor central para processamento e armazenamento. A solução inclui mecanismos de ingestão via API REST, persistência em banco de dados relacional, painel \textit{web} para visualização e regras de detecção de padrões como tentativas de autenticação malsucedidas e varreduras de portas, com notificação automática por Microsoft Teams e e-mail. O sistema foi operado em rede local controlada com múltiplos clientes Linux reais, demonstrando a viabilidade prática da integração entre monitoramento de infraestrutura e análise de eventos de segurança.

    \palavraschave{Monitoramento de Sistemas; Servidores Linux; Análise de Logs; Observabilidade; Arquitetura de Software}
\end{resumo}

\listadefiguras
\listadecodigos

\begin{siglas}
    \item[CPU]    \textit{Central Processing Unit} (Unidade Central de Processamento)
    \item[SRE]    \textit{Site Reliability Engineering} (Engenharia de Confiabilidade de Sites)
    \item[RAM]    \textit{Random Access Memory} (Memória de Acesso Aleatório)
    \item[I/O]    \textit{Input/Output} (Entrada/Saída)
    \item[R/W]    \textit{Read/Write} (Leitura/Escrita)
    \item[IOPS]   \textit{Input/Output Operations Per Second} (Operações de Entrada/Saída por Segundo)
    \item[JSON]   \textit{JavaScript Object Notation} (Notação de Objetos JavaScript)
    \item[JSONB]  \textit{JSON Binary} (Formato de armazenamento binário para JSON no PostgreSQL)
    \item[SNMP]   \textit{Simple Network Management Protocol} (Protocolo Simples de Gerenciamento de Redes)
    \item[MIB]    \textit{Management Information Base} (Base de Informações de Gerenciamento)
    \item[HDD]    \textit{Hard Disk Drive} (Unidade de Disco Rígido)
    \item[SSD]    \textit{Solid State Drive} (Unidade de Estado Sólido)
    \item[IDS]    \textit{Intrusion Detection System} (Sistema de Detecção de Intrusão)
    \item[SSH]    \textit{Secure Shell} (Terminal Seguro)
    \item[SYN]    \textit{Synchronize} (Pacote TCP de sincronização)
    \item[FIN]    \textit{Finish} (Pacote TCP de finalização)
    \item[LGPD]   Lei Geral de Proteção de Dados
    \item[API]    \textit{Application Programming Interface} (Interface de Programação de Aplicações)
    \item[XML]    \textit{eXtensible Markup Language} (Linguagem de Marcação Extensível)
    \item[HTTP]   \textit{Hypertext Transfer Protocol} (Protocolo de Transferência de Hipertexto)
    \item[TCP]    \textit{Transmission Control Protocol} (Protocolo de Controle de Transmissão)
\end{siglas}


% Lista de símbolos
%\begin{simbolos}
%    \item[$ \Sigma $]   Somatório
%    \item[$ \prod $]    Produtório
%    \item[$ \in $]      Pertence 
%    \item[$ \notin $]   Não pertence
%    \item[$ \mathcal{O}(n) $] Notação Big O 
%    \item[$ \forall $]  Para todo (quantificador universal)
%    \item[$ \exists $]  Existe (quantificador existencial)
%\end{simbolos}

\sumario

\textual

% -------------------------------------------------------------------------------------------------
% -------------------------------------------------------------------------------------------------
% 1. INTRODUÇÃO
% -------------------------------------------------------------------------------------------------
% -------------------------------------------------------------------------------------------------
\chapter{Introdução}
\label{cap:introducao}

Atualmente, em ambientes corporativos, servidores desempenham papel essencial na sustentação de aplicações, serviços internos e comunicação em rede. A indisponibilidade de recursos, falhas operacionais ou incidentes de segurança podem impactar diretamente a continuidade dos serviços, tornando necessária a adoção de mecanismos que permitam acompanhar o estado dos sistemas de forma contínua. Nesse contexto, o monitoramento de sistemas possibilita a identificação de falhas, degradações de desempenho e indisponibilidades a partir da coleta e análise de dados operacionais. De acordo com a \citeonline{fortinet_monitoring}, o monitoramento consiste na observação contínua da infraestrutura com o objetivo de detectar problemas e garantir a disponibilidade dos serviços.

Além disso, com o aumento da complexidade dos ambientes computacionais, observa-se a necessidade de ampliar a capacidade de análise desses sistemas, incorporando conceitos como observabilidade. Segundo \citeonline{ibm_observability}, a observabilidade permite compreender o comportamento dos sistemas a partir de dados como métricas e registros de eventos. Nesse cenário, os logs assumem papel relevante na identificação de atividades suspeitas e na análise de incidentes de segurança, sendo considerados componentes fundamentais em processos de monitoramento, conforme destacado por \citeonline{nist_sp800_92}.

Apesar da existência de diversas ferramentas no mercado, muitas soluções apresentam elevado nível de complexidade ou exigem infraestrutura robusta, o que dificulta sua adoção em ambientes de pequeno porte ou em contextos acadêmicos. Dessa forma, torna-se relevante investigar como os conceitos de monitoramento e análise de eventos podem ser aplicados na construção de soluções simplificadas e funcionais.

Diante desse contexto, este trabalho tem como problema de pesquisa a seguinte questão: como implementar um sistema capaz de monitorar métricas de infraestrutura e identificar eventos básicos de segurança em um ambiente de rede local?

O objetivo geral deste projeto é desenvolver um sistema de monitoramento para ambientes Linux, capaz de coletar métricas operacionais e identificar eventos de segurança a partir da análise de logs. Como objetivos específicos, destacam-se: (i) implementar mecanismos de coleta de dados nos dispositivos monitorados; (ii) desenvolver um servidor central para processamento e armazenamento das informações; (iii) identificar eventos de segurança com base em padrões conhecidos, como tentativas de autenticação malsucedidas e varreduras de rede; e (iv) disponibilizar uma interface para visualização dos dados e geração de alertas.

Para atingir esses objetivos, foi adotada uma abordagem baseada em arquitetura cliente-servidor, com coleta de dados realizada por agentes e envio das informações a um servidor central para processamento e armazenamento.

Este trabalho está organizado da seguinte forma: os capítulos iniciais apresentam a fundamentação teórica relacionada ao tema. Em seguida, são descritos o planejamento e a metodologia de desenvolvimento. Os capítulos finais detalham a implementação do sistema, os resultados de validação e as considerações acerca das limitações e trabalhos futuros.

% Corrigir última linha, para faze-la se adequa-la ao projeto atual usando as labels
% -------------------------------------------------------------------------------------------------
% -------------------------------------------------------------------------------------------------
% -------------------------------------------------------------------------------------------------

% -------------------------------------------------------------------------------------------------
% -------------------------------------------------------------------------------------------------
% 2. Monitoramento de Sistemas
% -------------------------------------------------------------------------------------------------
% -------------------------------------------------------------------------------------------------

\chapter{Monitoramento de Sistemas}
\label{cap:mon-sys}

Este capítulo dedica-se a explorar os fundamentos do monitoramento, iniciando pela sua definição formal e discutindo sua importância estratégica nas organizações modernas. Em seguida, são detalhadas as principais métricas utilizadas para aferir o  desempenho da infraestrutura, com ênfase na abordagem metodológica USE (Utilização, Saturação e Erros), amplamente adotada para diagnóstico de gargalos. Por fim, são abordados os mecanismos de alerta, essenciais para a transição de uma postura reativa para uma atuação proativa na gestão de incidentes.

\section{Definição}

O monitoramento de sistemas consiste no processo contínuo de coleta, exame e compreensão de dados sobre o funcionamento e o estado dos sistemas, assegurando que operem corretamente. Esse processo envolve a observação de métricas e eventos em tempo real, o que ajuda a identificar falhas, quedas de desempenho e comportamentos anormais.

De acordo com \citeonline{tanenbaum2015sistemas}, vigiar constantemente os componentes do sistema é crucial para administrar sistemas operacionais de maneira eficaz, especialmente no que se refere ao uso de CPU, memória e dispositivos de entrada e saída. Adicionalmente, pesquisas como a de \citeonline{aceto2019cloud} ressaltam que o monitoramento desempenha um papel fundamental em ambientes distribuídos, ao permitir a constante coleta de dados sobre o estado do sistema e auxiliar métodos de gestão e adequação flexível.

\section{Importância}

A importância do monitoramento de sistemas está diretamente relacionada à necessidade de garantir disponibilidade, desempenho e confiabilidade dos serviços de TI. Nos cenários atuais, nos quais as aplicações dependem fortemente da infraestrutura computacional, a ausência de monitoramento prejudica significativamente a agilidade para solucionar erros.

De acordo com a \citeonline{ibm2022server}, o monitoramento ajuda a identificar problemas antes que eles afetem os usuários finais, diminuindo o tempo de inatividade do sistema e aumentando a eficácia das operações. Conforme discutido por \citeonline{beyer2016site} no modelo de Engenharia de Confiabilidade do Google (SRE), o monitoramento contínuo possibilita acompanhar a condição do sistema em tempo real, oferecendo dados cruciais para o planejamento de capacidade e a melhoria do desempenho por meio do estudo de padrões ao longo do tempo.

Outro ponto relevante é o auxílio na tomada de decisão. A partir dos dados coletados, gestores podem identificar padrões de uso, prever demandas futuras e ajustar a infraestrutura de forma mais eficiente. Desse modo, o monitoramento torna-se mais do que uma ferramenta técnica, configurando-se como um recurso estratégico nas empresas, uma vez que o uso de dados analíticos contribui para decisões mais assertivas e competitivas \cite{davenport2006competing}.

\section{Métricas de Monitoramento}

As métricas são dados numéricos que servem para medir o desempenho dos recursos computacionais. Segundo \citeonline{gregg2020systems}, a análise de métricas como uso de CPU, memória e armazenamento é fundamental para identificar gargalos e compreender o comportamento de sistemas computacionais, o que ajuda a fazer diagnósticos melhores e otimizar de forma eficaz.

\subsection{CPU (Processador)}

O acompanhamento do uso do processador é crucial, pois está diretamente relacionado à capacidade de processamento do sistema. Essa métrica geralmente é expressa em percentual, indicando o quanto do recurso está sendo utilizado em determinado momento. Segundo a \citeonline{aws2023ec2}, altos níveis de utilização da CPU podem indicar sobrecarga, enquanto valores constantemente baixos podem apontar subutilização de recursos. Além disso, a análise dessa métrica ao longo do tempo permite identificar gargalos e necessidades de escalabilidade.

\subsection{Memória RAM}

A memória RAM atua como um espaço de trabalho temporário, guardando as informações que os programas em uso necessitam, sendo crucial para que os computadores operem sem engasgos. O monitoramento da memória é importante para prevenir situações como lentidão e travamentos. De acordo com \citeonline{tanenbaum2015sistemas}, a insuficiência de memória pode levar ao uso excessivo de memória virtual (\textit{swap}), o que impacta negativamente o desempenho do sistema. Métricas como memória utilizada, memória livre e uso de \textit{swap} são amplamente utilizadas para avaliar a eficiência do uso desse recurso.

\subsection{Disco (Armazenamento)}

O monitoramento de disco envolve a análise das operações de entrada e saída (E/S), englobando a gravação e o acesso a dados. Esse tipo de métrica é essencial para identificar gargalos relacionados ao armazenamento. Segundo \citeonline{barrett2016performance}, métricas como IOPS (operações de entrada/saída por segundo), latência e taxa de transferência são fundamentais para avaliar o desempenho do subsistema de armazenamento. Problemas de disco podem impactar diretamente o tempo de resposta das aplicações e a experiência do usuário.

\subsection{Abordagens de Análise de Métricas}

Uma abordagem amplamente utilizada na análise de métricas é o método USE (\textit{Utilization, Saturation, Errors}), proposto por Brendan Gregg. Esse método consiste na análise de três aspectos principais dos recursos computacionais:

\begin{itemize}
    \item \textbf{Utilização}: nível de uso do recurso ao longo do tempo;
    \item \textbf{Saturação}: presença de filas ou sinais de sobrecarga, indicando que o recurso está no limite de sua capacidade;
    \item \textbf{Erros}: falhas ocorridas durante a operação do sistema.
\end{itemize}

De acordo com \citeonline{gregg2013use}, essa abordagem é eficaz para a identificação sistemática de problemas de desempenho, pois permite uma análise estruturada e abrangente dos principais recursos do sistema.

\section{Alertas no Monitoramento}

Os alertas são mecanismos essenciais dentro do monitoramento de sistemas, responsáveis por notificar automaticamente administradores ou outros sistemas quando determinadas condições são atendidas, como a ultrapassagem de limites previamente definidos. A configuração de alertas possibilita uma resposta rápida a incidentes, reduzindo o impacto de falhas e melhorando a disponibilidade dos serviços \cite{googlecloud2023alerting}. Esses alertas podem ser configurados com base em limites fixos, padrões históricos ou técnicas de detecção de anomalias.

Sistemas modernos de monitoramento utilizam diferentes tipos de alertas, incluindo:
\begin{itemize}
    \item Alertas baseados em limiar (\textit{threshold});
    \item Alertas baseados em comportamento;
    \item Alertas baseados em anomalias.
\end{itemize}

A correta configuração de alertas é fundamental para evitar problemas como o excesso de notificações, conhecido como \textit{alert fatigue}, que pode reduzir a efetividade do monitoramento e comprometer a resposta a incidentes \cite{ibm2022itmonitoring}.
% -------------------------------------------------------------------------------------------------
% -------------------------------------------------------------------------------------------------
% -------------------------------------------------------------------------------------------------

% -------------------------------------------------------------------------------------------------
% -------------------------------------------------------------------------------------------------
% 3. Observabilidade
% -------------------------------------------------------------------------------------------------
% -------------------------------------------------------------------------------------------------

\chapter{Observabilidade}
\label{cap:observabilidade}

Diferente do monitoramento, que foca no estado do sistema, a observabilidade é uma propriedade intrínseca dele. Conforme \citeonline{majors2022observability}, isso possibilita que o desenvolvedor de software rastreie um problema (o sintoma) até a sua origem, sem necessitar de hipóteses iniciais.

\section{Diferença entre Monitoramento e Observabilidade}

Embora os termos sejam frequentemente utilizados como sinônimos, a literatura técnica moderna estabelece distinções fundamentais entre eles. Conforme as práticas de Engenharia de Confiabilidade de \citeonline{beyer2016site}, o monitoramento é o processo de coletar, agregar e analisar métricas em tempo real para entender o estado de um sistema em relação a um conjunto de indicadores pré-definidos. Essa ação se concentra no que está acontecendo, sendo muito importante para o funcionamento básico.

Por outro lado, a observabilidade, definida por Schwartz (2017), é o quanto é possível saber sobre o funcionamento interno de um sistema apenas observando seus sinais externos. Enquanto o monitoramento age quando algo já deu errado e se baseia em falhas já conhecidas, a observabilidade é proativa e antecipa problemas. Conforme explicam \citeonline{majors2022observability}, a observabilidade ajuda os engenheiros a investigar problemas que nunca apareceram antes (os chamados \textit{unknown unknowns}), permitindo que eles naveguem pelos dados de telemetria sem precisar criar alertas para cada possível erro.

\section{Os Três Pilares da Observabilidade}

A base de um sistema observável é construída sobre três categorias diferentes de dados de telemetria: métricas, logs e \textit{traces}. Como demonstrado por \citeonline{sigelman2010dapper}, o segredo para entender tudo o que acontece em sistemas complexos está em juntar esses dados. As métricas mostram se algo está errado usando números gerais, enquanto os logs guardam um histórico completo e detalhado de tudo que aconteceu. E, para completar, os \textit{traces} mostram o caminho que cada solicitação faz, desde o início até o fim, ajudando a encontrar exatamente onde está o problema ou a lentidão em um sistema complicado \cite{majors2022observability}.

\subsection{Métricas}

As métricas são representações numéricas de dados agregados ao longo de intervalos de tempo. Segundo a especificação do \citeonline{opentelemetry2023spec}, elas são otimizadas para armazenamento e consulta rápida, sendo o pilar ideal para o monitoramento de saúde do sistema e acionamento de alertas.

\subsection{Registros de Eventos (Logs)}

Diferente das métricas, os logs consistem em anotações textuais distintas de ocorrências em um instante determinado. Para \citeonline{majors2022observability}, o log formatado (como no padrão JSON) apresenta maior utilidade, já que possibilita a busca por características particulares dentro da comunicação por meio de aplicativos de análise. Nos aprofundaremos mais nos logs no capítulo~\ref{cap:logs}.

\subsection{Rastreamento Distribuído (Distributed Tracing)}

O rastreamento é o pilar que conecta os pontos em uma arquitetura de microsserviços, mantendo o registro do caminho de cada solicitação enquanto ela passa por diferentes partes da estrutura. A base técnica para este pilar é o artigo ``Dapper'' \cite{sigelman2010dapper}, que introduziu os conceitos de \textit{span} e \textit{trace ID}. Além disso, \citeonline{cederberg2021observability} ressalta que, em ambientes com microsserviços, o rastreamento distribuído é essencial para enxergar como os serviços autônomos se relacionam no tempo, permitindo a identificação de problemas interligados que passariam despercebidos em registros isolados. Cabe ressaltar, contudo, que o rastreamento distribuído é mais relevante em arquiteturas de microsserviços; o sistema implementado neste trabalho adota uma arquitetura centralizada e contempla apenas métricas e \textit{logs}, de modo que esse pilar permanece fora do escopo da solução.

% -------------------------------------------------------------------------------------------------
% -------------------------------------------------------------------------------------------------
% -------------------------------------------------------------------------------------------------

% -------------------------------------------------------------------------------------------------
% -------------------------------------------------------------------------------------------------
% 4. Coleta de Dados
% -------------------------------------------------------------------------------------------------
% -------------------------------------------------------------------------------------------------
\chapter{Coleta de Dados de um Endpoint}
\label{cap:data-collect}

A coleta de dados de um \textit{endpoint}\footnote{Endpoint refere-se a um ponto final de comunicação em uma rede. O termo pode designar tanto um dispositivo participante, como um servidor ou estação de trabalho, quanto um ponto específico de acesso a um serviço, como uma URL em uma API.} e o modelo de comunicação são etapas fundamentais na arquitetura de sistemas, pois permitem a obtenção contínua de informações e a transferência de dados entre cliente e servidor, que será mais explicado no capítulo~\ref{cap:archteture}. Esses dados representam atualizações e mudanças de estado, sendo essenciais para o gerenciamento em tempo real das aplicações.

\section{Coleta de Dados em Sistemas de Monitoramento}

A coleta de dados é uma etapa fundamental em sistemas de monitoramento, especialmente, mas não se limitando a ambientes Linux, pois permite a obtenção contínua de informações sobre o desempenho e o estado dos recursos computacionais. Esses dados incluem métricas como uso de CPU, memória, disco e tráfego de rede, sendo essenciais para o gerenciamento da infraestrutura.

No contexto de monitoramento de redes, o uso de agentes é uma abordagem amplamente adotada. Conforme descrito por \citeonline[p.~14]{santos2019monitoramento}, “o agente é responsável por coletar as informações do dispositivo monitorado e disponibilizá-las ao sistema gerenciador”. Esse modelo permite que os dados sejam capturados diretamente no sistema operacional e enviados para análise centralizada.

Um dos principais protocolos utilizados nesse processo é o SNMP (\textit{Simple Network Management Protocol}), que possibilita a coleta estruturada de informações por meio de objetos organizados em uma base denominada MIB (\textit{Management Information Base}). De acordo com \citeonline{case2002snmp}, o SNMP foi desenvolvido com o objetivo de facilitar o gerenciamento de redes, permitindo a leitura e manipulação de informações de dispositivos de forma padronizada.

Além disso, a arquitetura de monitoramento envolve três componentes principais: o agente (capítulo~\ref{cap:agents}), o dispositivo monitorado (\textit{endpoint}) e o gerenciador(servidor). O agente coleta os dados localmente, enquanto o gerenciador é responsável por centralizar e analisar essas informações. Em sistemas práticos, como os que utilizam ferramentas de monitoramento, essa estrutura é essencial para garantir eficiência e escalabilidade.

A coleta de dados pode ocorrer por diferentes métodos, sendo os principais o \textit{polling} e o orientado a eventos, que veremos mais adiante. Conforme apontado por \citeonline[p.~6]{souza2020nagios}, “o monitoramento pode ser realizado por verificações periódicas (\textit{polling}) ou por notificações automáticas enviadas pelos dispositivos”. Essa combinação permite maior confiabilidade na detecção de falhas.

\section{Coleta Periódica de Dados e suas Utilidades}

A coleta periódica de dados consiste na obtenção de informações em intervalos regulares de tempo, sendo uma prática essencial em sistemas de monitoramento. Essa abordagem permite a criação de históricos que auxiliam na análise do comportamento dos sistemas ao longo do tempo.

De acordo com \citeonline[p.~22]{ferreira2020monitoramento}, “a coleta contínua de dados possibilita o acompanhamento da evolução dos recursos computacionais e a identificação de padrões de uso”. Essa característica é fundamental para ambientes que exigem alta disponibilidade e desempenho.

Uma das principais vantagens da coleta periódica é a capacidade de detectar anomalias e prever falhas antes que causem impactos significativos. Nesse sentido, a análise de dados históricos desempenha papel essencial. Conforme destacado por \citeonline[p.~35]{lima2021planejamento}, “a análise de dados coletados ao longo do tempo permite identificar comportamentos fora do padrão e agir de forma preventiva”.

Além disso, a coleta periódica contribui para o planejamento de capacidade, pois fornece informações sobre o crescimento do uso de recursos. Segundo \citeonline[p.~24]{ferreira2020monitoramento}, esses dados “auxiliam no dimensionamento adequado da infraestrutura e na tomada de decisões estratégicas”.

Outro aspecto importante refere-se à definição da frequência de coleta. Intervalos muito curtos podem gerar sobrecarga no sistema, enquanto intervalos muito longos podem comprometer a precisão das análises. Conforme observado por \citeonline[p.~18]{santos2019monitoramento}, “a escolha da periodicidade de coleta deve equilibrar desempenho e qualidade das informações obtidas”.

Por fim, a coleta periódica permite a geração de relatórios, auditorias e análises históricas, sendo essencial para a gestão eficiente dos sistemas computacionais.

\section{Armazenamento de Dados e suas Estruturas}

O armazenamento de dados consiste na retenção de informações digitais em meios físicos, sendo uma prática essencial em sistemas computacionais para as operações de software. Essa abordagem permite a criação de um ecossistema sólido que auxilia na consulta, análise e uso das informações pelas aplicações.

De acordo com \citeonline{susnjara_storage}, “armazenamento de dados refere-se a mídias magnéticas, ópticas ou mecânicas que registram e preservam informações digitais para operações contínuas ou futuras”. Essa característica física, baseada em memórias RAM, HDDs e SSDs, é fundamental para ambientes que exigem alta disponibilidade e desempenho.

Uma das principais decisões no armazenamento é a escolha do modelo de banco de dados, que organiza os registros para que não causem impactos significativos na aplicação. Nesse sentido, os bancos relacionais desempenham papel essencial. Conforme destacado pela \citeonline{aws2026database}, “os bancos de dados relacionais armazenam dados em formato de tabela, consistindo em linhas e colunas”.

Além disso, o formato físico de armazenamento contribui para a transferência e a confiabilidade da infraestrutura, como ocorre no armazenamento em blocos. Segundo \citeonline{susnjara_storage}, nesse formato os dados “são armazenados como peças separadas, cada uma com um identificador único”.

Outro aspecto importante refere-se à escolha da manutenção da infraestrutura, dividida entre modelos autogerenciados e serviços na nuvem. Modelos autogerenciados exigem grande esforço técnico operacional, enquanto a nuvem traz conveniência. Conforme observado pela \citeonline{aws2026database}, “os benefícios de usar um serviço gerenciado incluem a redução do tempo gasto em gerenciamento e manutenção da infraestrutura e maior segurança”.

Por fim, o correto armazenamento físico e lógico permite a geração de valor, consultas estruturadas e integração de serviços, sendo essencial para a gestão eficiente das plataformas e dos sistemas corporativos.

\section{Tipo de requisição de Coleta de dados}

No contexto de verificação de atualizações, o uso de \textit{polling} é uma abordagem amplamente adotada. Conforme \citeonline{rios2024polling}, o \textit{polling} “refere-se a uma técnica de comunicação onde um cliente solicita regularmente, ou em intervalos regulares, informações de um servidor”. Esse modelo permite que os dados sejam verificados de maneira ativa pela aplicação cliente, por meio de requisições cíclicas.

Um dos principais modelos utilizados para otimizar esse processo é a arquitetura orientada a eventos (\textit{event-driven}), que possibilita uma comunicação estruturada reagindo a mudanças de estado. De acordo com a \citeonline{aws2025eda}, a arquitetura baseada em eventos “usa eventos para acionar e comunicar entre serviços desacoplados”, permitindo a manipulação de informações de forma assíncrona e independente.

Além disso, a arquitetura orientada a eventos envolve três componentes principais: produtores, \textit{brokers}\footnote{\textit{Broker} é o componente intermediário responsável por receber, filtrar e distribuir mensagens entre produtores e consumidores. O termo é frequentemente traduzido como ``roteador de mensagens'', porém essa tradução pode gerar confusão com roteadores de rede (dispositivos físicos que encaminham pacotes IP); por isso, o termo em inglês é preferível no contexto de arquiteturas orientadas a eventos.} e consumidores. O produtor emite os eventos; o \textit{broker} é responsável por centralizar, filtrar e distribuir essas mensagens aos consumidores. Em sistemas práticos modernos, essa estrutura é essencial para garantir eficiência, tolerância a falhas e escalabilidade.

\subsection{Exemplo didático: diferença entre polling e event-driven}

Para compreender melhor a diferença entre os dois modelos, imagine o controle de acesso a dois edifícios diferentes:

\textbf{Modelo polling, prédio com detecção facial automática:} O sistema de câmeras e sensores funciona continuamente, verificando o rosto de quem se aproxima a cada fração de segundo. Mesmo quando não há ninguém, o sistema segue “perguntando” ao servidor: “Alguém quer entrar?”, centenas de vezes por minuto. Isso gera muitas verificações desnecessárias (chamadas vazias), mas o acesso é liberado imediatamente quando uma pessoa autorizada é detectada.

\textbf{Modelo event-driven, prédio com porteiro eletrônico:} O sistema fica ocioso até que alguém toque a campainha (o evento). Só então o porteiro (o roteador) é acionado, verifica a autorização e libera a entrada. Nenhuma verificação ocorre enquanto não há um evento. Isso elimina chamadas vazias, mas a resposta depende da ocorrência do evento.

Em resumo: no \textit{polling} o cliente pergunta repetidamente (ativo); no \textit{event-driven} o servidor avisa quando algo muda (reativo). A escolha depende do cenário: \textit{polling} é simples e fácil de implementar, enquanto \textit{event-driven} é mais eficiente e escalável para sistemas com muitos eventos esporádicos.

A comunicação de dados pode ocorrer por diferentes métodos, sendo os principais o \textit{polling} e o modelo orientado a eventos. Conforme apontado pela \citeonline{microsoft2026eda}, no modelo orientado a eventos, “os eventos são entregues quase em tempo real, para que os consumidores possam responder imediatamente aos eventos assim que ocorrerem”. Essa combinação elimina as chamadas vazias e garante maior agilidade no fluxo de dados da rede.

% -------------------------------------------------------------------------------------------------
% -------------------------------------------------------------------------------------------------
% -------------------------------------------------------------------------------------------------

% -------------------------------------------------------------------------------------------------
% -------------------------------------------------------------------------------------------------
% 5. Dashboards e visualização de Dados
% -------------------------------------------------------------------------------------------------
% -------------------------------------------------------------------------------------------------

\chapter{Visualização e Dashboards}
\label{cap:dash}

Os \textit{dashboards}, também conhecidos como painéis de informação, são ferramentas fundamentais na gestão moderna, pois permitem organizar, analisar e visualizar dados de forma clara, interativa e consolidada em um único ambiente \cite{ufms2024}. Sua principal função é transformar dados brutos em \textit{insights} estratégicos, facilitando a compreensão de informações complexas e apoiando a tomada de decisão de maneira rápida, precisa e baseada em dados concretos \cite{rolim2020}.

Essas ferramentas oferecem uma visão abrangente e atualizada ao centralizar dados provenientes de diferentes fontes, possibilitando o acompanhamento do desempenho e dos resultados organizacionais em tempo real \cite{ufms2024}. Por meio de elementos visuais como gráficos de barras, gráficos de setores, mapas interativos e cartões de métricas, os \textit{dashboards} tornam a interpretação das informações mais intuitiva, permitindo identificar padrões, tendências e anomalias com maior facilidade \cite{campos2023}. Além disso, a interatividade, com filtros, menus e navegação por abas, possibilita que o usuário explore diferentes recortes dos dados, enquanto recursos como o \textit{drill-down} permitem aprofundar a análise e obter informações mais detalhadas \cite{campos2023}.

Outro aspecto relevante é a capacidade de sintetizar grandes volumes de dados, reduzindo a complexidade e evitando sobrecarga de informações, o que é especialmente importante em contextos com dados heterogêneos \cite{rolim2020}. A integração de análises multidimensionais e geoespaciais também amplia a compreensão dos dados, permitindo visualizar relações espaciais e temporais de forma clara \cite{ufms2024}.

Os \textit{dashboards} ainda contribuem diretamente para a eficiência operacional, ao possibilitar o monitoramento contínuo e em tempo real de indicadores, facilitando a identificação imediata de problemas e a adoção de ações corretivas \cite{ufms2024}. Além disso, promovem vantagem competitiva ao otimizar o tempo de análise, reduzir custos e melhorar a qualidade dos serviços, substituindo relatórios tradicionais por visualizações dinâmicas e interativas \cite{rolim2020}.

Por fim, destacam-se pela acessibilidade, permitindo que diferentes perfis de usuários, mesmo sem grande conhecimento técnico, consigam interpretar dados e extrair \textit{insights} relevantes para embasar decisões estratégicas \cite{ufms2024}. Dessa forma, os \textit{dashboards} atuam como uma ponte entre a coleta de dados e a ação, simplificando a análise de informações complexas e tornando a gestão mais eficiente, ágil e orientada por dados \cite{campos2023}.

% -------------------------------------------------------------------------------------------------
% -------------------------------------------------------------------------------------------------
% -------------------------------------------------------------------------------------------------

% -------------------------------------------------------------------------------------------------
% -------------------------------------------------------------------------------------------------
% 6. Agentes de Monitoramento
% -------------------------------------------------------------------------------------------------
% -------------------------------------------------------------------------------------------------
\chapter{Agentes de Monitoramento}
\label{cap:agents}

Conforme falado no capítulo ~\ref{cap:data-collect} agentes de monitoramento são programas instalados nos dispositivos com a finalidade de coletar informações sobre o estado e o desempenho dos sistemas computacionais. Esses agentes atuam diretamente no sistema operacional, capturando métricas como uso de CPU, memória, disco e processos em execução, sendo essenciais para o acompanhamento da infraestrutura de TI.

De acordo com \citeonline[p.~13]{santos2019monitoramento}, “os agentes de monitoramento são programas executados nos dispositivos que coletam informações e as disponibilizam para sistemas gerenciadores”. Dessa forma, eles funcionam como intermediários entre os dispositivos monitorados e os sistemas centrais, garantindo que os dados sejam coletados de forma organizada e confiável.

O funcionamento desses agentes baseia-se na coleta contínua ou sob demanda de dados, que são posteriormente enviados para um servidor central para análise. Conforme destacado por \citeonline[p.~21]{ferreira2020monitoramento}, “os agentes realizam a coleta contínua de dados do sistema e os encaminham para análise”, permitindo o acompanhamento do comportamento dos recursos ao longo do tempo e a identificação de possíveis falhas.

Além disso, os agentes podem operar em diferentes modos, como ativo e passivo. No modo passivo, respondem a requisições do sistema gerenciador (\textit{polling}), enquanto no modo ativo enviam informações automaticamente quando ocorrem eventos relevantes. Segundo \citeonline[p.~5]{souza2020nagios}, “os sistemas de monitoramento podem operar com verificações periódicas ou envio automático de alertas”, o que aumenta a eficiência na detecção de problemas.

Por fim, destaca-se que os agentes frequentemente utilizam protocolos padronizados, como o SNMP\footnote{O SNMP (\textit{Simple Network Management Protocol}) é um protocolo padrão para coleta e gerenciamento de informações sobre dispositivos em redes IP. Os dados são organizados em uma estrutura hierárquica chamada MIB (\textit{Management Information Base}), que define quais variáveis o agente pode expor ao sistema gerenciador, como uso de CPU, interfaces de rede e contadores de erros.}, para comunicação com os sistemas de monitoramento. Conforme estabelecido por \citeonline{case2002snmp}, esses protocolos permitem a padronização da coleta de dados, facilitando a interoperabilidade entre diferentes dispositivos e tornando o monitoramento mais escalável e eficiente.

\section{Modelo Agente - Servidor}


Enquanto no modelo cliente-servidor o cliente não compartilha nenhum de seus recursos com o servidor, é ele quem inicia a comunicação solicitando alguma função ao servidor que fica aguardando as requisições, com o agente a iniciativa da comunicação parte dele próprio: o agente pode reportar dados de forma periódica (\textit{push} em intervalos regulares) ou orientada a eventos, contactando o servidor quando há algo relevante a reportar.

\citeonline[p.~116]{wooldridge1995} identificam seis propriedades do agente: a \textbf{autonomia}, em que os agentes operam sem a intervenção direta de humanos ou outros e têm algum tipo de controle sobre suas ações e estado interno; a \textbf{habilidade social}, na qual os agentes interagem com outros agentes ou possivelmente humanos por meio de algum tipo de linguagem de comunicação; a \textbf{reatividade}, onde os agentes percebem seu ambiente, que pode ser o mundo físico, um usuário, uma coleção de outros agentes ou a Internet, e respondem em tempo hábil às mudanças que ocorrem nele; e a \textbf{proatividade}, em que os agentes não agem simplesmente em resposta ao seu ambiente, sendo capazes de exibir comportamento direcionado a objetivos e tomar a iniciativa.

\section{Ferramentas de Monitoramento: Zabbix e Nagios}

Diversas ferramentas utilizam agentes de monitoramento para coletar dados e acompanhar o funcionamento dos sistemas, sendo o Zabbix e o Nagios duas das soluções mais utilizadas tanto no meio acadêmico quanto no mercado. Essas ferramentas exemplificam, na prática, a aplicação dos conceitos de monitoramento em ambientes reais.

O Nagios é uma ferramenta consolidada que permite monitorar \textit{hosts}, serviços e dispositivos de rede por meio de agentes e \textit{plugins}. Conforme afirmam \citeonline[p.~7]{souza2020nagios}, “o Nagios permite o monitoramento de \textit{hosts} e serviços, utilizando \textit{plugins} e agentes para coleta de dados em diferentes ambientes”, o que demonstra sua flexibilidade e ampla aplicabilidade.

Uma característica importante do Nagios é a possibilidade de operar com verificações ativas e passivas, permitindo adaptar o monitoramento conforme as necessidades do ambiente. Essa abordagem possibilita tanto a coleta periódica de dados quanto o envio de alertas automáticos, contribuindo para a rápida identificação de falhas.

O Zabbix, por sua vez, destaca-se pelo monitoramento em tempo real e pela capacidade de lidar com grandes volumes de dados. Ele utiliza agentes instalados nos dispositivos para coletar informações detalhadas sobre o desempenho do sistema. Segundo \citeonline[p.~25]{ferreira2020monitoramento}, o Zabbix “permite monitoramento em tempo real de diversos parâmetros, com suporte a coleta distribuída de dados”, o que o torna adequado para ambientes mais complexos.

Dessa forma, tanto o Zabbix quanto o Nagios demonstram a importância dos agentes de monitoramento na coleta e análise de dados. Essas ferramentas permitem a geração de alertas, relatórios e análises de desempenho, contribuindo para a gestão eficiente dos sistemas e garantindo maior confiabilidade na infraestrutura de TI.

% -------------------------------------------------------------------------------------------------
% -------------------------------------------------------------------------------------------------
% -------------------------------------------------------------------------------------------------

% -------------------------------------------------------------------------------------------------
% -------------------------------------------------------------------------------------------------
% 7. Logs e Analise de Eventos
% -------------------------------------------------------------------------------------------------
% -------------------------------------------------------------------------------------------------
\chapter{Logs e Análise de Eventos}
\label{cap:logs}

Os logs são registros de eventos gerados por sistemas, aplicações e dispositivos ao longo de seu funcionamento, nos quais cada entrada representa uma ocorrência específica. Segundo \citeonline{nist_sp800_92}, “um log é um registro dos eventos que ocorrem nos sistemas e redes de uma organização”. De acordo com a \citeonline{ibm_log_analysis}, esses registros constituem dados essenciais para as operações de TI, pois fornecem informações relevantes sobre desempenho e segurança.

A importância dos logs está relacionada à visibilidade que proporcionam e à capacidade de reconstruir ações dentro do ambiente. \citeonline{todd2017logging} destaca que operar uma rede sem logs limita significativamente a compreensão dos eventos ocorridos, dificultando a identificação de atividades maliciosas. \citeonline{gupta2012logging} complementa que a análise de eventos em tempo quase real é fundamental para minimizar impactos de incidentes e atender a requisitos de conformidade. Para uma gestão eficiente, a infraestrutura de \textit{logging}\footnote{\textit{Logging} refere-se ao processo sistemático de registrar eventos, ações e estados de um sistema em arquivos de log. O termo, derivado do inglês \textit{log} (registro), engloba tanto a geração dos registros pelos sistemas quanto o seu armazenamento, centralização e análise posterior.} deve ser organizada em camadas funcionais.

Conforme \citeonline{nist_sp800_92}, esse processo envolve a geração de logs por \textit{hosts} e dispositivos, o armazenamento e a análise por servidores responsáveis pela centralização dos dados, e o monitoramento por meio de interfaces que permitem a visualização e interpretação das informações. Os registros podem ser classificados de acordo com sua finalidade.

A \citeonline{ibm_log_analysis} destaca os logs de acesso, que registram o comportamento de usuários; os logs de erro, que documentam falhas e incidentes; e os logs de eventos, que representam mudanças de estado no sistema. Independentemente da origem, a padronização é essencial para garantir interoperabilidade.

Segundo \citeonline{todd2017logging}, logs bem estruturados devem conter, no mínimo, informações como data, hora, fuso horário, origem, nível de severidade e descrição do evento. A análise de logs tem como objetivo transformar dados brutos em informações úteis para a tomada de decisão. Esse processo envolve etapas como a normalização (\textit{parsing}), que consiste na conversão de dados não estruturados em formatos organizados, facilitando sua interpretação.

De acordo com o \citeonline{gupta2012logging}, essa etapa permite a correlação de eventos provenientes de múltiplas fontes, contribuindo para a identificação de ataques mais complexos. Além disso, \citeonline{todd2017logging} destaca a importância da redução de dados desnecessários, visando otimizar armazenamento sem comprometer o valor analítico dos registros.

Como exemplo prático, considere o seguinte registro de um dispositivo de segurança:


\begin{codigo}[htb]
\legenda[exemplo_log]{Exemplo de um log gerado por um Firewall}
\begin{lstlisting}
    date=2026-04-05 time=14:32:10 devname=FW-01 type=traffic 
    subtype=forward srcip=192.168.1.45 dstip=8.8.8.8 action=deny 
    protocol=TCP port=443 msg="Blocked outbound connection"
\end{lstlisting}
\fonteautor
\end{codigo}

\begin{verbatim}
\end{verbatim}

Nesse caso, é possível identificar que o dispositivo ``FW-01'' bloqueou uma tentativa de comunicação externa realizada a partir de um endereço IP interno, utilizando o protocolo TCP na porta 443. A análise estruturada desses campos permite compreender o evento registrado e identificar possíveis comportamentos anômalos, sendo fundamental para atividades de monitoramento e segurança.

% -------------------------------------------------------------------------------------------------
% -------------------------------------------------------------------------------------------------
% -------------------------------------------------------------------------------------------------

% -------------------------------------------------------------------------------------------------
% -------------------------------------------------------------------------------------------------
% 8. Detecção de Eventos de Segurança
% -------------------------------------------------------------------------------------------------
% -------------------------------------------------------------------------------------------------

\chapter{Detecção de Eventos de Segurança}
\label{cap:detection}

A coleta massiva de dados de monitoramento e logs, discutida nos capítulos anteriores, só atinge seu pleno potencial quando associada a mecanismos eficazes de análise e interpretação. No âmbito da segurança da informação, essa análise converte registros brutos em inteligência acionável contra ameaças cibernéticas. Este capítulo concentra-se na detecção de eventos de segurança a partir da observação de comportamentos anômalos na rede e nos sistemas. Inicialmente, apresentam-se os princípios gerais que norteiam a identificação de atividades  maliciosas por meio de Sistemas de Detecção de Intrusão (IDS). Na sequência, são examinados dois vetores de ataque comuns em ambientes Linux e amplamente documentados na literatura: os ataques de força bruta, que visam comprometer credenciais de acesso, e as varreduras de portas (\textit{port scanning}), utilizadas para reconhecimento de infraestrutura. A compreensão desses padrões fundamenta a definição das regras de correlação de eventos que integrarão a proposta do sistema de monitoramento.

\section{detecção de eventos de segurança e como são detectados?}

A detecção de eventos de segurança tem como base a análise de \textit{logs}, que são registros detalhados de tudo o que acontece em um sistema, aplicação ou rede. Conforme discutido no capítulo anterior, registros são essenciais para que os profissionais de TI possam entender o comportamento de seus ambientes e identificar possíveis ameaças. Como explica a \citeonline{ibm_log_analysis}, ferramentas avançadas de análise de \textit{logs} podem até automatizar a detecção de atividades suspeitas, alertando os gestores quando algo anormal ocorre.

Para detectar intrusões, os Sistemas de Detecção de Intrusão (IDS) monitoram o tráfego de rede em busca de padrões maliciosos. Conforme \citeonline{patel2016}, os sistemas baseados em regras (também chamados de detecção por uso indevido) comparam o tráfego com um banco de assinaturas de ataques conhecidos, sendo muito eficazes contra ameaças já catalogadas. Os autores afirmam que “a detecção baseada em regras é muito eficaz contra ataques conhecidos”. No entanto, \citeonline{sbseg2023} alertam que “os métodos tradicionais para detectar varreduras de portas são limitados porque se baseiam em regras estáticas e no conhecimento prévio da estrutura da rede”. Para superar essa limitação, \citeonline{salgueiro2010} propõem o uso de uma linguagem declarativa específica para descrever assinaturas de intrusão que podem se espalhar por vários pacotes de rede, aumentando a capacidade de detecção.

\section{Ataques de força bruta}

O ataque de força bruta é uma tentativa repetitiva de adivinhar senhas e credenciais de acesso por meio de tentativa e erro. Como descreve a \citeonline{ibm_bruteforce}, “um ataque de força bruta é um tipo de ciberataque no qual \textit{hackers} tentam obter acesso não autorizado a uma conta ou a dados criptografados por tentativa e erro, testando várias credenciais de \textit{login}”. O serviço SSH (\textit{Secure Shell}), amplamente utilizado para acesso remoto a servidores, é um dos principais alvos desse tipo de ataque. De acordo com a \citeonline{startupdefense_ssh}, “um aumento acentuado nas tentativas de \textit{login} com falha é um dos indicadores mais claros de um ataque de força bruta”. Além disso, \citeonline{park2021network} afirmam que a maioria dos \textit{logs} gerados por acessos não autorizados em roteadores de internet é causada por ataques de força bruta SSH, o que permite coletar os endereços IP dos atacantes.

Os riscos de um ataque bem-sucedido são altos. A \citeonline{fortinet_bruteforce} destaca que “um ataque de força bruta bem-sucedido pode resultar em acesso não autorizado imediato, permitindo que os atacantes se façam passar pelo usuário [e] roubem dados sensíveis”. Uma vez dentro do sistema, o invasor pode usar os recursos do servidor para fins maliciosos, como participar de \textit{botnets} ou lançar novos ataques. Por isso, monitorar \textit{logs} de falhas de \textit{login} é uma prática de segurança fundamental.

\section{Ataques de footprinting na rede (port scan)}

Antes de tentar invadir um sistema, o atacante geralmente realiza um reconhecimento da rede, técnica conhecida como \textit{footprinting}. A principal ferramenta desse reconhecimento é a varredura de portas (\textit{port scan}). Segundo \citeonline{arxiv2301_13581}, “a varredura de portas é o processo de tentar se conectar a várias portas de rede em um dispositivo para determinar quais estão abertas e quais serviços estão sendo executados”. O autor complementa que o primeiro passo de qualquer ataque cibernético é universalmente entendido como o reconhecimento, e que esse reconhecimento quase sempre inclui algum tipo de varredura de portas.

Na prática, o invasor envia pacotes especiais para cada porta do alvo e analisa as respostas. \citeonline{patel2016} explicam que, por meio dessa técnica, o atacante identifica vulnerabilidades e fraquezas nos pontos de entrada da máquina. Existem diferentes tipos de varredura, como a SYN scan, a FIN scan e a Null scan, que tentam se passar por tráfego normal para evitar a detecção. \citeonline{sbseg2023} apresentam um método que detecta varreduras de portas mesmo sem conhecimento prévio da estrutura da rede, o que representa um avanço em relação às abordagens tradicionais.

\section{A importância de saber detectar esses ataques}

Detectar varreduras de portas e ataques de força bruta logo no início é essencial para evitar danos maiores. \citeonline{arxiv2301_13581} ressalta que “ao identificar tentativas de varredura de portas, os profissionais de cibersegurança podem tomar medidas proativas para proteger os sistemas e redes antes que um invasor tenha a chance de explorar qualquer vulnerabilidade”. Em outras palavras, a detecção precoce transforma \textit{logs} passivos em uma defesa ativa. Conforme afirmam \citeonline{sbseg2023}, “a varredura de portas é um primeiro passo em diferentes vetores de ataque. Portanto, é essencial detectar essas varreduras para limitar os seus impactos”.

Além da proteção imediata, a análise de \textit{logs} auxilia no cumprimento de regulamentações como a Lei Geral de Proteção de Dados (LGPD) e padrões do setor financeiro, que exigem rastreabilidade e registros de acesso. Como destaca \citeonline{nist_sp800_92}, a análise regular de \textit{logs} ajuda os profissionais de TI a entender melhor como seus sistemas estão funcionando, melhorar o desempenho e aumentar a segurança. Em resumo, saber detectar esses ataques não é apenas uma questão técnica, é uma necessidade estratégica para qualquer organização que queira proteger seus dados e sua reputação.
% -------------------------------------------------------------------------------------------------
% -------------------------------------------------------------------------------------------------
% -------------------------------------------------------------------------------------------------


% -------------------------------------------------------------------------------------------------
% -------------------------------------------------------------------------------------------------
% 9. Comunicação e Arquitetura
% -------------------------------------------------------------------------------------------------
% -------------------------------------------------------------------------------------------------

\chapter{Comunicação e Arquitetura}
\label{cap:archteture}

O modelo cliente-servidor estrutura a comunicação entre programas em dois papéis: o cliente é o responsável por iniciar a requisição e o servidor é responsável por processá-la e retornar uma resposta. Nesse modelo, a comunicação entre cliente e servidor ocorre por meio de uma API que atua como intermediária: quando o cliente quer algo, ele faz uma requisição seguindo as regras definidas pela API, o servidor recebe essa requisição, a processa e devolve uma resposta. Essas regras determinam como os dados devem ser enviados, em qual formato e quais operações são permitidas, garantindo que os dois lados se comuniquem de forma padronizada e organizada.

A Figura~\ref{fig:cliente-servidor} ilustra esse fluxo de comunicação: o cliente envia uma solicitação, a API intermedeia o processamento no servidor e a resposta retorna ao cliente por meio dessa mesma camada.

\begin{figure}[htb]
\legenda[fig:cliente-servidor]{Fluxo de comunicação no modelo cliente-servidor intermediado por uma API}
\fig{width=0.95\textwidth}{../img/cliente-servidor.png}
\fonte{Os autores (2026).}
\end{figure}

De acordo com \citeonline[p.~24]{oliveira2016}, os principais componentes desse modelo são os servidores, responsáveis por oferecer um conjunto de serviços, como servidores de impressão, de arquivos e com \textit{web services}; os clientes, que usam os serviços disponibilizados nos servidores, sendo tipicamente subsistemas executados independentemente que, embora informados sobre os servidores disponíveis, desconhecem a existência de outros clientes, destacando-se que seus processos são separados; e uma rede que viabilize o acesso por parte dos clientes aos serviços, componente este desnecessário apenas quando cliente e servidor encontram-se na mesma máquina. Ressalta-se, no entanto, que sistemas que fazem uso de uma arquitetura cliente-servidor utilizam um sistema distribuído.
    
JSON é um formato de texto para organizar dados. Não pertence a nenhuma linguagem de programação (define só estrutura), então qualquer sistema consegue ler o que outro escreveu. Segundo a \citeonline{ibm_json}, os dados ficam em pares de chave e valor, o que facilita tanto a leitura por máquina quanto a interpretação de quem está desenvolvendo. Essa padronização também ajuda na troca de informações entre sistemas diferentes, já que todos seguem a mesma lógica de estrutura.

Sua estrutura leve e compacta o torna mais eficiente que o XML, que carrega mais marcações e ocupa mais espaço para representar a mesma informação. Essa diferença de peso aparece especialmente em sistemas que fazem muitas requisições, onde cada \textit{kilobyte} a menos no \textit{payload} conta. Em aplicações modernas, como APIs e sistemas distribuídos, essa economia de dados contribui diretamente para melhor desempenho e menor consumo de recursos.

HTTP\footnote{O HTTPS (\textit{Hypertext Transfer Protocol Secure}) é a variante segura do HTTP, que adiciona uma camada de cifragem TLS (\textit{Transport Layer Security}), garantindo confidencialidade e integridade na transmissão dos dados. O sistema de monitoramento implementado neste trabalho utiliza HTTP puro, limitação discutida no Capítulo~\ref{cap:limitacoes}.} é o protocolo que move dados na \textit{web}. O modelo é simples: cliente pede, servidor responde. Esse funcionamento é a base da \textit{web}, como destaca a \citeonline{mdn_http}. Cada troca é independente, o que facilita escalar sistemas sem aumentar a complexidade. Isso permite que aplicações atendam um grande número de requisições sem precisar manter informações sobre interações anteriores.

Os métodos GET, POST, PUT e DELETE indicam o tipo de operação. Os cabeçalhos carregam informações sobre formato dos dados e autenticação. Os códigos de status (200 para sucesso, 404 para não encontrado) dizem ao cliente o que aconteceu com cada pedido, tornando o diagnóstico de erros mais direto. De acordo com a \citeonline{rfc9110}, esses elementos fazem parte da padronização do protocolo HTTP, garantindo uma comunicação consistente entre sistemas.

APIs conectam sistemas diferentes sem expor como funcionam por dentro. Elas recebem requisições, processam e respondem, mantendo a complexidade encapsulada e deixando os sistemas mais fáceis de manter. Isso permite que diferentes partes de um sistema evoluam de forma independente, sem impactar diretamente outras áreas.

Na prática, APIs definem como as requisições e respostas devem ser estruturadas: definem o que pode ser pedido, em qual formato e o que será devolvido. Como descrito na \citeonline{mdn_api}, esse modelo permite que partes diferentes de um sistema sejam desenvolvidas de forma independente. Um aplicativo móvel pode puxar dados de um \textit{backend} sem saber nada sobre o banco de dados, o que facilita a integração entre diferentes tecnologias e plataformas.

O Zabbix mostra bem como essas três tecnologias funcionam juntas. De acordo com a documentação oficial, a ferramenta usa APIs baseadas em HTTP para comunicação, enquanto os dados são transmitidos em JSON \cite{zabbix}. Isso permite coleta e análise em tempo real, com integração a outros sistemas sem exigir adaptações complexas em cada ponta. Na prática, isso possibilita monitorar servidores, redes e aplicações de forma centralizada, além de automatizar alertas e acompanhar métricas importantes em tempo real. O resultado é um fluxo de monitoramento contínuo, onde as informações trafegam em formato leve e são acessíveis por qualquer sistema que fale HTTP.

% -------------------------------------------------------------------------------------------------
% -------------------------------------------------------------------------------------------------
% -------------------------------------------------------------------------------------------------

% -------------------------------------------------------------------------------------------------
% -------------------------------------------------------------------------------------------------
% 10. Planejamento e arquitetura do Sistema proposto
% -------------------------------------------------------------------------------------------------
% -------------------------------------------------------------------------------------------------

\chapter{Aplicação dos Conceitos no Planejamento do Sistema de Monitoramento}
\label{cap:plan}

A proposta de um sistema de monitoramento, conforme delineada no Capítulo~\ref{cap:introducao}, pressupõe a capacidade de coletar métricas operacionais, armazenar \textit{logs}, gerar alertas e identificar eventos relevantes a partir da análise desses dados. Entretanto, a implementação dessas funcionalidades não é trivial e envolve uma série de decisões arquiteturais que dependem diretamente dos conceitos discutidos na fundamentação teórica.

Neste capítulo, apresenta-se o planejamento da solução proposta, evidenciando como os conceitos abordados nos Capítulos~\ref{cap:mon-sys} a \ref{cap:archteture} fundamentam as escolhas realizadas. Mais do que descrever a implementação, busca-se demonstrar que decisões como modelo de coleta, arquitetura de comunicação, estrutura de armazenamento e mecanismos de detecção de eventos só se tornam compreensíveis à luz da teoria previamente apresentada.

\section{Visão Geral da Solução Proposta}

O sistema adota uma arquitetura centralizada, na qual um servidor único é responsável por receber, processar e armazenar dados provenientes de múltiplos clientes monitorados. Essa escolha está diretamente relacionada ao modelo agente‑servidor discutido no Capítulo~\ref{cap:agents}, no qual agentes distribuídos coletam informações localmente e as encaminham a um gerenciador central.

Apesar de arquiteturas distribuídas oferecerem maior escalabilidade e tolerância a falhas, optou-se por um modelo centralizado devido ao escopo controlado do projeto e à redução da complexidade de implementação. Essa decisão, no entanto, implica um ponto único de falha, o que será discutido posteriormente como limitação da solução.

Os clientes executam agentes de monitoramento responsáveis por coletar dados do sistema operacional e enviá-los ao servidor. Essa abordagem segue o princípio de agentes autônomos descrito por \citeonline{wooldridge1995}, nos quais entidades são capazes de perceber o ambiente e agir de forma proativa.

No que se refere à comunicação, adotou-se o modelo \textit{push}, no qual os próprios agentes enviam os dados ao servidor em intervalos regulares. Essa escolha está alinhada com os conceitos apresentados no Capítulo~\ref{cap:data-collect}, especialmente no que diz respeito ao balanceamento entre frequência de coleta e impacto na rede. Diferentemente do modelo \textit{polling}, o \textit{push} reduz requisições desnecessárias, porém pode gerar picos de carga no servidor caso múltiplos agentes enviem dados simultaneamente.

\section{Arquitetura Física e Lógica}

A arquitetura física do sistema é composta por um servidor central e seis clientes, todos conectados a um \textit{switch} não gerenciável em topologia estrela. A rede utiliza a faixa \texttt{10.10.10.0/26} (máscara \texttt{255.255.255.192}), permitindo até 62 endereços IP utilizáveis, com atribuição dinâmica realizada por meio de um serviço DHCP executado no próprio servidor.

Essa configuração simplifica o gerenciamento da rede, mas limita a capacidade de controle sobre o tráfego, uma vez que o \textit{switch} utilizado não permite segmentação ou priorização de pacotes, o que pode impactar o desempenho em cenários de maior carga.

A Figura~\ref{fig:topologia-pkt} ilustra a disposição dos componentes na rede local. O servidor central (\texttt{10.10.10.1}) executa os serviços de \textit{backend}, banco de dados, DHCP e \textit{proxy} reverso. Os clientes recebem endereços dinamicamente pelo DHCP na faixa \texttt{10.10.10.10}--\texttt{10.10.10.50} e executam o agente coletor. Todos os dispositivos são interligados por um \textit{switch} não gerenciável, caracterizando uma topologia estrela.

\begin{figure}[htb]
\legenda[fig:topologia-pkt]{Topologia lógica do sistema de monitoramento}
\fig{width=1.00\textwidth}{../img/topologia.png}
\fonte{Os autores (2026).}
\end{figure}


No plano lógico, o servidor é estruturado utilizando contêineres Docker. A containerização consiste no empacotamento de uma aplicação e suas dependências em unidades isoladas e portáteis, garantindo que o ambiente de execução seja reproduzível independentemente da máquina hospedeira. Essa abordagem reduz conflitos de dependência e facilita a implantação dos componentes do sistema, embora não elimine a necessidade de gerenciamento de recursos, especialmente em ambientes com limitação de CPU e memória.

Os principais componentes do servidor são:

\begin{itemize}
    \item \textbf{Nginx}: atua como \textit{proxy} reverso, recebendo requisições HTTP e encaminhando‑as ao \textit{backend}. Essa camada permite futura adoção de HTTPS e centraliza o controle de acesso.
    \item \textbf{Backend (Python/Flask)}: responsável por expor a API REST, validar requisições e processar os dados recebidos. A autenticação é realizada por meio de uma chave de API dinâmica (\textit{API key}), gerada pelo painel e armazenada no banco de dados.
    \item \textbf{PostgreSQL}: utilizado para armazenamento de métricas e \textit{logs}. A escolha por um banco relacional se deve à necessidade de consultas estruturadas, embora isso possa limitar a escalabilidade em cenários com alto volume de escrita contínua.
    \item \textbf{Frontend (React 19 + Vite + TailwindCSS + react-chartjs-2 + react-gauge-component)}: responsável pela visualização dos dados, consumindo informações via API REST.
\end{itemize}

\section{Coleta de Dados e Agentes de Monitoramento}

A coleta de dados implementada pelos agentes segue dois modelos distintos: coleta inicial (\textit{discovery}) e coleta contínua de métricas.

A coleta inicial tem como objetivo estabelecer uma linha de base do sistema monitorado, incluindo informações de \textit{hardware}, sistema operacional e \textit{software} instalado. Essa etapa é essencial para contextualizar as métricas coletadas posteriormente.

Já a coleta contínua ocorre em intervalos de aproximadamente 5 segundos, capturando métricas como uso de CPU, memória, disco e rede. A escolha desse intervalo reflete um compromisso entre granularidade dos dados e sobrecarga no sistema, conforme discutido no Capítulo~\ref{cap:mon-sys}.

A seleção das métricas baseia-se no método USE (\textit{Utilization, Saturation, Errors}), proposto por \citeonline{gregg2013use}, que orienta a análise de desempenho a partir da utilização de recursos, níveis de saturação e ocorrência de erros.

Além das métricas, os agentes coletam \textit{logs} do sistema operacional, permitindo a análise de eventos e a identificação de comportamentos suspeitos. A centralização desses \textit{logs} segue as recomendações de \citeonline{nist_sp800_92}, que destaca sua importância para auditoria e investigação de incidentes.

Os dados são enviados ao servidor via requisições HTTP POST em formato JSON. Embora essa abordagem seja simples e amplamente suportada, ela não garante confidencialidade ou integridade dos dados em sua forma básica, sendo necessária a adoção de HTTPS em versões futuras.

\section{Armazenamento e Estruturação dos Dados}

O sistema utiliza o PostgreSQL como banco de dados principal, estruturando as informações em tabelas relacionais para \textit{hosts}, métricas, \textit{logs} e alertas.

A utilização do tipo JSONB permite armazenar dados semiestruturados, o que oferece flexibilidade na evolução do sistema sem necessidade de alterações frequentes no \textit{schema}.

Entretanto, o uso de um banco relacional único pode se tornar um gargalo em cenários com alta taxa de ingestão de dados, especialmente devido à natureza contínua das métricas coletadas.

Para mitigar o crescimento descontrolado do banco, foi definida e implementada uma política de retenção de dados de 7 dias, executada automaticamente por um agendador de tarefas integrado ao \textit{backend}. A limpeza ocorre diariamente às 03:00 UTC, eliminando registros de métricas e \textit{logs} com mais de sete dias. Esse período equilibra a necessidade de análise histórica de curto prazo, suficiente para identificar tendências semanais e investigar incidentes recentes, com a limitação de armazenamento disponível no ambiente acadêmico.

\section{Processamento de Eventos e Alertas}

O processamento dos dados ocorre por meio de regras definidas no \textit{backend}, que avaliam tanto métricas quanto \textit{logs}.

No contexto de monitoramento de infraestrutura, são utilizados limiares (\textit{thresholds}) para geração de alertas, como uso elevado de CPU ou baixo espaço em disco. Essa abordagem é simples e eficaz, mas depende de calibração adequada para evitar excesso de notificações.

No que se refere à segurança, a análise baseia-se em padrões conhecidos identificados nos \textit{logs}, como múltiplas tentativas de autenticação malsucedidas (indicando possível ataque de força bruta) ou conexões em múltiplas portas (indicando varredura de portas). Conforme discutido no Capítulo~\ref{cap:detection}, sistemas baseados exclusivamente em regras são eficazes contra ameaças conhecidas, mas podem falhar na identificação de comportamentos inéditos. A opção por essa abordagem simplificada é compatível com o objetivo de prova de conceito do presente trabalho, reduzindo a complexidade de implementação sem comprometer a validação dos princípios de detecção estudados.

\section{Visualização e Dashboard}

O \textit{dashboard web} apresenta os dados coletados por meio de gráficos e tabelas, permitindo a visualização do estado atual do sistema e a análise de tendências.

As principais funcionalidades incluem:

\begin{itemize}
    \item Visualização de métricas em séries temporais;
    \item Exibição de \textit{logs} recentes;
    \item Listagem de alertas ativos;
    \item \textit{Status} dos \textit{hosts} monitorados.
\end{itemize}

Embora a interface permita a análise operacional do ambiente, sua eficácia depende diretamente da qualidade dos dados coletados e das regras definidas no \textit{backend}.

As limitações identificadas na solução são detalhadas em conjunto com os trabalhos futuros no Capítulo~\ref{cap:limitacoes}.

O planejamento aqui descrito evidencia que cada componente do sistema de monitoramento encontra respaldo direto nos capítulos teóricos precedentes. A compreensão dos conceitos de métricas, \textit{logs}, observabilidade, arquitetura cliente‑servidor e segurança da informação mostrou-se imprescindível para a tomada de decisões arquiteturais conscientes e para a construção de uma solução funcional e didaticamente relevante. Os capítulos seguintes detalham a metodologia de desenvolvimento adotada, a implementação concreta do sistema e os resultados obtidos com sua operação em ambiente controlado.

% -------------------------------------------------------------------------------------------------
% -------------------------------------------------------------------------------------------------
% -------------------------------------------------------------------------------------------------

% =================================================================================================
% 11. METODOLOGIA DE DESENVOLVIMENTO
% =================================================================================================
\chapter{Metodologia de Desenvolvimento}
\label{cap:metodologia}

Este capítulo descreve a abordagem adotada pelo grupo para organizar e conduzir o
desenvolvimento do sistema, as ferramentas utilizadas e o cronograma de execução.

\section{Abordagem e Organização do Projeto}

O desenvolvimento foi conduzido por uma equipe de sete integrantes, com divisão de
responsabilidades orientada por componente. A infraestrutura de rede e do ambiente
(contêineres Docker, DNS e DHCP) e a coordenação geral do projeto ficaram a cargo de
Lucas Toterol Rodrigues, que também desenvolveu o agente coletor, com o apoio de Caio
Federico Esquivel Lovera Arze, responsável ainda pela redação técnica do documento. O
desenvolvimento do \textit{backend} (API REST e lógica de detecção) coube a Douglas
Clayton da Silva e Fernando Bach; o do \textit{frontend}, a Denyse Panza Clemente Ferreira
e Larissa Quirino dos Santos; e a modelagem do banco de dados, a Beatriz Pimentel de Mello.

A metodologia adotada foi o desenvolvimento iterativo e incremental. O grupo organizou
o trabalho em ciclos semanais de entrega, nos quais cada membro contribuía com
incrementos ao repositório compartilhado. As reuniões de equipe ocorriam, em geral, aos
fins de semana, nem sempre com a presença de todos os integrantes.

O controle de versão foi realizado integralmente por meio do Git, com repositório
hospedado no GitHub. O histórico de \textit{commits} registra 202 revisões ao longo
do período de desenvolvimento, refletindo a evolução incremental de todos os
componentes do sistema.

\section{Cronograma e Fases}
\label{sec:cronograma}

O desenvolvimento ocorreu entre março e junho de 2026, totalizando 202 \textit{commits}
no repositório. Por se tratar de um processo iterativo e incremental, as fases a seguir
não são estritamente sequenciais: houve sobreposição entre elas, em especial entre a
implementação dos componentes e o refinamento subsequente. O recorte abaixo, reconstituído
a partir do histórico de \textit{commits}, organiza o trabalho pela ênfase predominante em
cada período.

O trabalho foi norteado por um cronograma inicial de sete semanas (21 de abril a 5 de
junho de 2026), organizado por entregas semanais e por componente. A execução acompanhou
esse planejamento em linhas gerais, porém com desvios: o arranque efetivo da implementação
ocorreu no início de maio, cerca de uma semana após o previsto, e os refinamentos do
\textit{frontend} estenderam-se até meados de junho, ultrapassando o congelamento planejado
do MVP (5 de junho). Em decorrência desses ajustes, a apresentação final do projeto foi
realizada em 26 de junho de 2026.

\begin{itemize}
    \item \textbf{Fase 1: Pesquisa e Especificação}: levantamento bibliográfico,
          definição da arquitetura e elaboração da fundamentação teórica (Capítulos 1 a
          10). Essa etapa corresponde à fase inicial do projeto e antecede o versionamento
          contínuo do código no repositório atual.
    \item \textbf{Fase 2: Infraestrutura e Prova de Conceito}: configuração do
          servidor Ubuntu, rede local \texttt{10.10.10.0/26}, DHCP e DNS, e primeiros
          contêineres Docker (março e abril de 2026).
    \item \textbf{Fase 3: Implementação dos Componentes}: desenvolvimento paralelo do
          agente coletor, do \textit{backend} e do \textit{frontend}, com integração
          progressiva (do final de abril a maio de 2026, período de maior concentração de
          \textit{commits}). O refinamento da interface do \textit{frontend} prosseguiu
          até meados de junho.
    \item \textbf{Fase 4: Refinamento e Endurecimento}: migração para autenticação por
          chave de API dinâmica, remoção de tabelas obsoletas, redução da superfície de
          ataque e adição do agendador de limpeza automática (início de junho de 2026).
    \item \textbf{Fase 5: Validação}: execução dos cenários de teste de detecção
          (força bruta e varredura de portas) em rede local controlada, realizada ao final
          do ciclo de desenvolvimento. A consolidação dos resultados é apresentada no
          Capítulo~\ref{cap:resultados}.
\end{itemize}

\section{Ferramentas e Ambiente de Desenvolvimento}
\label{sec:ferramentas}

As ferramentas utilizadas no desenvolvimento estão distribuídas por camada do sistema:

\begin{itemize}
    \item \textbf{Controle de versão}: Git e GitHub.
    \item \textbf{Ambiente de execução}: Docker e Docker Compose para contêinerização
          dos serviços do servidor.
    \item \textbf{Linguagem do \textit{backend} e agente}: Python 3.12, com as
          bibliotecas Flask 3.0.3 (API REST), SQLAlchemy 2.0.30 (mapeamento
          objeto-relacional), APScheduler 3.10.4 (agendamento de tarefas), psutil
          (métricas do sistema operacional) e requests (comunicação HTTP).
    \item \textbf{Empacotamento do agente}: PyInstaller (\textit{flag}
          \texttt{--onefile}), gerando um único binário executável para distribuição
          nos clientes Linux.
    \item \textbf{Serviço de inicialização do agente}: systemd
          (\texttt{linux-agent.service}).
    \item \textbf{Linguagem do \textit{frontend}}: JavaScript com React 19, Vite
          (\textit{build tool}), TailwindCSS, react-router-dom 7 (roteamento),
          react-chartjs-2 (gráficos de linha) e react-gauge-component (medidores).
    \item \textbf{Banco de dados}: PostgreSQL 15.
    \item \textbf{Proxy reverso}: Nginx.
    \item \textbf{Infraestrutura de rede}: Ubuntu Server como hospedeiro físico, com
          \texttt{isc-dhcp-server} e \texttt{dnsmasq} instalados diretamente no sistema
          operacional do servidor.
    \item \textbf{Ambiente de desenvolvimento (IDE)}: Visual Studio Code.
    \item \textbf{Gestão e documentação}: Jira (gestão de tarefas) e Confluence
          (documentação técnica colaborativa).
\end{itemize}

% =================================================================================================
% 12. IMPLEMENTAÇÃO DO SISTEMA
% =================================================================================================
\chapter{Implementação do Sistema}
\label{cap:implementacao}

Este capítulo descreve a implementação concreta do sistema de monitoramento, detalhando
cada componente, suas tecnologias e decisões de projeto.

\section{Visão Geral da Arquitetura Implementada}
\label{sec:arq-impl}

A arquitetura implementada segue o modelo planejado no Capítulo~\ref{cap:plan}. O
servidor central executa quatro contêineres Docker orquestrados pelo Docker Compose:
\textbf{nginx} (\textit{proxy} reverso, única porta externamente exposta: 80),
\textbf{backend} (API REST em Flask), \textbf{frontend} (SPA React) e
\textbf{postgres} (PostgreSQL 15). Os clientes executam o agente coletor como
serviço systemd. Toda a comunicação entre agente e servidor ocorre via HTTP,
endereçada ao nome \texttt{api.monitoramento.lan}, resolvido pelo DNS local.

A Figura~\ref{fig:arquitetura} apresenta o diagrama de arquitetura do sistema, com os
componentes e suas conexões; a Figura~\ref{fig:fluxograma} ilustra o fluxo de dados desde
a coleta no agente até a geração de alertas e o envio de notificações.

\begin{figure}[htb]
\legenda[fig:arquitetura]{Diagrama de arquitetura do sistema implementado}
\fig{width=1.00\textwidth}{../img/arquitetura.png}
\fonte{Os autores (2026).}
\end{figure}

\begin{figure}[htb]
\legenda[fig:fluxograma]{Fluxograma do fluxo de dados da aplicação}
\fig{width=1.00\textwidth}{../img/fluxograma.png}
\fonte{Os autores (2026).}
\end{figure}

\section{Infraestrutura de Rede e Servidores}
\label{sec:infra}

O ambiente de rede é composto por um servidor Ubuntu Server e múltiplos clientes
Linux, conectados em topologia estrela por meio de um \textit{switch} não gerenciável.
A rede utiliza a sub-rede \texttt{10.10.10.0/26} (máscara \texttt{255.255.255.192},
\textit{broadcast} \texttt{10.10.10.63}), com 62 endereços utilizáveis.

O servidor hospedeiro executa dois serviços nativos, fora dos contêineres: o
\texttt{isc-dhcp-server}, que distribui endereços na faixa
\texttt{10.10.10.10--10.10.10.50}, e o \texttt{dnsmasq}, que resolve os nomes
\texttt{api.monitoramento.lan} e \texttt{painel.monitoramento.lan} para o endereço
do servidor (\texttt{10.10.10.1}).

\subsection{Contêineres e Exposição de Portas}

A única porta exposta externamente é a 80. As portas internas 5432 (PostgreSQL),
5000 (\textit{backend} Flask) e 5173 (Vite/\textit{frontend}) são acessíveis apenas
dentro da rede interna do Docker Compose, sem mapeamento para o hospedeiro. Essa
decisão de endurecimento reduz a superfície de ataque ao mínimo necessário.

\subsection{Proxy Reverso (Nginx)}

O Nginx atua como ponto de entrada único e roteia as requisições por nome de
domínio (\textit{server\_name}), e não por prefixo de caminho: as requisições
destinadas a \texttt{api.monitoramento.lan} são encaminhadas ao contêiner
\textit{backend}, e as destinadas a \texttt{painel.monitoramento.lan}, ao contêiner
\textit{frontend}. Ambos os nomes resolvem para o endereço do servidor por meio do
DNS local (\texttt{dnsmasq}). O servidor responde exclusivamente em HTTP, sem
terminação TLS, limitação discutida no Capítulo~\ref{cap:limitacoes}.

\section{Agente Coletor}
\label{sec:agente}

O agente é um processo Python empacotado como binário único pelo PyInstaller e
instalado nos clientes como serviço systemd (\texttt{linux-agent.service}). Sua
lógica principal é um laço de 5 segundos que executa, em ordem: \textit{flush} da
fila de \textit{retry}, envio de \textit{heartbeat}, coleta e envio de métricas,
coleta de linhas novas do \texttt{auth.log} e envio de detecção de
\textit{port scan} (quando detectada).

A instalação do agente nos clientes monitorados é automatizada pelo script \texttt{install.sh}, que compila o binário, cria o serviço systemd e o habilita. As Figuras~\ref{fig:install-1} e~\ref{fig:install-2} apresentam, respectivamente, o início do processo de instalação e sua conclusão, com o serviço \texttt{linux-agent.service} registrado e iniciado no host.

\begin{figure}[htb]
\legenda[fig:install-1]{Início da instalação do agente coletor pelo script \texttt{install.sh}}
\fig{width=1.00\textwidth}{../img/instalacao-1.png}
\fonte{Os autores (2026).}
\end{figure}

\begin{figure}[htb]
\legenda[fig:install-2]{Conclusão da instalação, com o serviço systemd habilitado e iniciado}
\fig{width=1.00\textwidth}{../img/instalacao-2.png}
\fonte{Os autores (2026).}
\end{figure}

\subsection{Discovery: Físico versus Virtual}
\label{subsec:discovery}

Antes de iniciar o laço principal, o agente realiza uma coleta única de inventário de
\textit{hardware} e \textit{software} (\textit{discovery}). A distinção entre máquina
física e virtual é feita pela presença do campo \texttt{Hypervisor Vendor} na saída do
comando \texttt{lscpu}: se presente, a coleta de dados de \textit{hardware} físico
(módulos de memória via \texttt{dmidecode}, discos via \texttt{smartctl}) é omitida,
coletando-se apenas os dados acessíveis ao sistema operacional convidado.

Os módulos de \textit{discovery} estão organizados por componente em subdiretórios de
\texttt{Agent/src/agent/discovery/}: \texttt{cpu\_discovery}, \texttt{disk\_discovery},
\texttt{mem\_discovery}, \texttt{network\_discovery}, \texttt{motherboard\_discovery},
\texttt{system\_discovery} e \texttt{tools\_discovery}. Cada um contém variantes
\texttt{\_physical.py} e \texttt{\_virtual.py}, selecionadas em tempo de execução.

Os dados de \textit{discovery} são enviados ao \textit{backend} via
\texttt{POST /api/discovery} e persistidos nas tabelas \texttt{host}, \texttt{agents}
e \texttt{host\_discovery}.

\subsection{Coleta Contínua de Métricas}
\label{subsec:metricas-agente}

A coleta de métricas é centralizada em \texttt{coleta/collector.py}, que agrega as
saídas dos subcoletores \texttt{cpu\_coleta}, \texttt{mem\_coleta},
\texttt{disk\_coleta}, \texttt{network\_coleta} e \texttt{process\_coleta}, todos
baseados na biblioteca \texttt{psutil}.

As métricas enviadas incluem: percentual de uso de CPU, uso e disponibilidade de
memória (em megabytes e percentual), uso e disponibilidade de disco, operações de
entrada/saída por segundo (IOPS) e taxa de transferência em bytes por segundo,
tráfego de rede e os 15 processos de maior consumo por CPU, memória, disco e
conexões. As métricas de IOPS e rede são calculadas como deltas entre ciclos
consecutivos; na primeira execução, esses valores são omitidos por ausência de
referência anterior.

\subsection{Coleta de Logs}
\label{subsec:logs-agente}

O módulo \texttt{coleta/logs\_coleta/logs.py} realiza leitura incremental do arquivo
\texttt{/var/log/auth.log}. A posição de leitura é rastreada por \textit{byte offset}
e número de \textit{inode}: na primeira execução, as últimas 100 linhas são enviadas;
nas execuções subsequentes, apenas as linhas novas desde o último ciclo. A rotação do
arquivo é detectada pela mudança de \textit{inode}, garantindo continuidade da coleta
sem perda de eventos.

\subsection{Detecção de Port Scan com tcpdump}
\label{subsec:portscan-agente}

A detecção de varredura de portas é realizada no próprio cliente monitorado, por meio
do subcomando \texttt{tcpdump} invocado como subprocesso. Dois \textit{threads} operam
em paralelo: o de captura lê continuamente pacotes SYN sem ACK; o de avaliação
(intervalo de 2 segundos) mantém uma janela deslizante de 60 segundos e conta o número
de portas destino distintas por IP de origem. Um IP é classificado como fonte de
varredura quando atinge o limiar de 10 portas distintas nessa janela.

Quando há fontes detectadas, o agente envia \texttt{POST /api/security/portscan} com
a lista de IPs. Ciclos sem detecção não geram tráfego de rede.

A ausência do \texttt{tcpdump} é tratada com degradação graciosa: o módulo emite
aviso no \textit{log} e desabilita a detecção silenciosamente. Para contornar uma
limitação do PyInstaller ao invocar subprocessos, que sobrescreve
\texttt{LD\_LIBRARY\_PATH}, o agente a restaura explicitamente antes de chamar
\texttt{tcpdump}.

O Código~\ref{cod:tcpdump-filtro} apresenta o filtro BPF (\textit{Berkeley Packet
Filter}) utilizado na chamada ao \texttt{tcpdump}. Ele captura exclusivamente
pacotes TCP que têm o bit SYN ativo e o bit ACK inativo, padrão dos pacotes de
abertura de conexão sem confirmação, predominante em ferramentas de varredura de
portas.

\begin{codigo}[htb]
\legenda[cod:tcpdump-filtro]{Filtro BPF para captura de pacotes TCP SYN no agente}
\begin{lstlisting}
tcpdump -n -l -i any \
    "tcp[tcpflags] & tcp-syn != 0 and tcp[tcpflags] & tcp-ack == 0"
\end{lstlisting}
\fonte{Adaptado de \texttt{Agent/src/agent/coleta/connections\_coleta/connections.py} (Os autores, 2026).}
\end{codigo}

\section{Backend: API REST}
\label{sec:backend}

O \textit{backend} é uma aplicação Flask organizada em \textit{blueprints}. Na
inicialização, executa 9 funções \texttt{garantir\_schema\_*()} que realizam migrações
idempotentes de esquema (\texttt{ALTER TABLE ... ADD COLUMN IF NOT EXISTS}), permitindo
que contêineres recriados a partir de volumes antigos sejam atualizados sem intervenção
manual. Um agendador APScheduler (\textit{BackgroundScheduler}) é iniciado na mesma
ocasião, programado para executar a limpeza do banco diariamente às 03:00 UTC.

Na inicialização do contêiner, o \textit{backend} também envia notificações Teams para
todos os alertas não resolvidos existentes no banco, sinalizando pendências ao operador.
Esse comportamento é protegido por uma \textit{flag} de arquivo para evitar repetição
em recarregamentos a quente do Flask.

\subsection{Endpoints Principais}
\label{subsec:endpoints}

Os \textit{endpoints} estão registrados em módulos de rota dedicados:

\begin{itemize}
    \item \textbf{Autenticação e configurações}: \texttt{POST /api/auth/login},
          \texttt{GET/POST /api/settings/apikey},
          \texttt{PATCH /api/settings/thresholds},
          \texttt{PATCH /api/settings/notifications},
          \texttt{GET/POST/DELETE /api/settings/email-recipients}.
    \item \textbf{Ingestão (protegidos por \texttt{X-API-Key})}:
          \texttt{POST /api/discovery}, \texttt{POST /api/metrics},
          \texttt{POST /api/logs}, \texttt{POST /api/security/portscan},
          \texttt{POST /api/heartbeat}.
    \item \textbf{Consulta (protegidos)}: \texttt{GET /api/hosts},
          \texttt{GET /api/metrics}, \texttt{GET /api/logs},
          \texttt{GET /api/alerts} (com filtros de \texttt{status} e
          \texttt{host\_id}).
    \item \textbf{Gestão}: \texttt{PATCH /api/alerts/<id>/resolve},
          \texttt{POST /api/maintenance/cleanup}.
    \item \textbf{Diagnóstico (abertos)}: \texttt{GET /health},
          \texttt{GET /api/status}.
\end{itemize}

\subsection{Normalização e Persistência}
\label{subsec:persistencia}

Os dados de \textit{log} são normalizados pelo módulo
\texttt{utils/parsers.py}, que aplica expressões regulares para identificar e
estruturar eventos do \texttt{auth.log}. São suportados seis tipos de evento:
autenticação SSH (\texttt{parse\_ssh\_log}), falha de autenticação PAM
(\texttt{parse\_pam\_auth\_failure}), sessão PAM (\texttt{parse\_pam\_session}),
sessão \textit{logind} (\texttt{parse\_logind\_session}), desconexão SSH
(\texttt{parse\_ssh\_disconnect}) e comandos sudo (\texttt{parse\_sudo}). Os dados
estruturados são persistidos em JSONB no campo \texttt{parsed\_data} da tabela
\texttt{logs}.

A persistência é mediada pelo ORM SQLAlchemy, com \textit{queries} exclusivamente
parametrizadas, eliminando a classe de vulnerabilidades de injeção SQL.

\subsection{Autenticação por Chave de API}
\label{subsec:autenticacao}

A autenticação é baseada no cabeçalho HTTP \texttt{X-API-Key}. A chave é gerada
pelo painel via \texttt{POST /api/settings/apikey/generate}, protegido pela senha
do painel (\texttt{PANEL\_PASSWORD}), utilizando \texttt{secrets.token\_hex(32)},
que produz uma cadeia hexadecimal de 64 caracteres (256 bits de entropia). A chave é
persistida na tabela \texttt{app\_settings} e lida pelos agentes a partir do arquivo
\texttt{/etc/monitor-agent/env} (permissões \texttt{600}).

O \textit{frontend} armazena a chave no \texttt{localStorage} do navegador e a
inclui em todas as requisições autenticadas. Esse mecanismo substituiu o
\textit{token} estático com cabeçalho \texttt{Authorization: Bearer} utilizado em
versões anteriores do sistema.

A Figura~\ref{fig:apikey} apresenta a tela de gerenciamento da chave de API no painel, na qual o operador pode gerar uma nova chave diretamente no servidor e configurá-la no navegador para autenticar as requisições.

\begin{figure}[htb]
\legenda[fig:apikey]{Tela de gerenciamento da chave de API no painel}
\fig{width=1.00\textwidth}{../img/tela-config-2-apiKEY.png}
\fonte{Os autores (2026).}
\end{figure}

\section{Detecção de Eventos e Geração de Alertas}
\label{sec:deteccao}

A lógica de detecção está centralizada em \texttt{utils/detection.py} e é acionada
pelas rotas de ingestão de métricas, \textit{logs} e \textit{port scan}.

\subsection{Detecção de Força Bruta SSH}
\label{subsec:bruteforce}

Ao receber uma linha de \textit{log} de autenticação com resultado \texttt{failed},
o \textit{backend} invoca \texttt{check\_brute\_force()}, que consulta a tabela
\texttt{logs} contando registros com \texttt{log\_type = 'auth'},
\texttt{status = 'failed'} e o mesmo IP de origem dentro de uma janela temporal de
60 segundos. Se o contador atingir ou superar 5 ocorrências, um alerta do tipo
\texttt{brute\_force} é criado, a menos que já exista um alerta ativo não resolvido
para o mesmo par (host, IP), evitando duplicações.

O Código~\ref{cod:bruteforce-query} apresenta o SQL equivalente executado para a
contagem de falhas. A consulta utiliza operadores JSONB nativos do PostgreSQL
(\texttt{->>}), aproveitando os índices parciais definidos sobre o campo
\texttt{parsed\_data} da tabela \texttt{logs}.

\begin{codigo}[htb]
\legenda[cod:bruteforce-query]{SQL de contagem de falhas SSH para detecção de força bruta}
\begin{lstlisting}
SELECT COUNT(*) FROM logs
WHERE host_id      = :host_id
  AND log_type     = 'auth'
  AND timestamp   >= :inicio_janela        -- utcnow - 60 s
  AND parsed_data->>'status'    = 'failed'
  AND parsed_data->>'ip_origem' = :ip_origem
\end{lstlisting}
\fonte{Adaptado de \texttt{Web/BackEnd/app/utils/parsers.py}, função \texttt{count\_failed\_logins()} (Os autores, 2026).}
\end{codigo}

\subsection{Detecção de Port Scan}
\label{subsec:portscan-backend}

O \textit{endpoint} \texttt{POST /api/security/portscan} recebe a lista de IPs
classificados como fontes de varredura pelo agente. Para cada IP,
\texttt{check\_port\_scan()} verifica se já existe um alerta ativo com
\textit{timestamp} nos últimos 2 minutos para o mesmo par (host, IP), o chamado \textit{cooldown}, evitando alertas repetidos em varreduras prolongadas.

Os dados de \textit{port scan} não são persistidos em tabela própria: o
\textit{endpoint} é estritamente sinalizador. A decisão de não persistir foi tomada
após a remoção da tabela \texttt{active\_connections} (junho de 2026), cujo custo de
manutenção superava o valor analítico dos dados armazenados.

\subsection{Alertas de Recurso (CPU, Memória e Disco)}
\label{subsec:resource-alerts}

A cada recebimento de métricas, o \textit{backend} chama
\texttt{check\_resource\_alert()} para CPU, memória e disco. O limiar de disparo é
lido da tabela \texttt{app\_settings} (padrão: 80\% caso não configurado). Um alerta
é gerado somente quando o valor excede o limiar e não existe alerta ativo do mesmo
tipo para o mesmo \textit{host}. Os limiares são ajustáveis em tempo de execução
pelo painel (\texttt{PATCH /api/settings/thresholds}).

\subsection{Notificações: Teams e E-mail}
\label{subsec:notificacoes}

Para cada alerta gerado, o \textit{backend} verifica os \textit{toggles}
\texttt{notify\_teams} e \texttt{notify\_email} na tabela \texttt{app\_settings}.
Se habilitadas, as funções \texttt{enviar\_alerta\_teams()} e
\texttt{enviar\_alerta\_email()} são invocadas.

As notificações via Microsoft Teams utilizam um \textit{webhook} do Power Automate,
com \textit{timeout} de 10 segundos e falha silenciosa. As notificações por e-mail
são enviadas via Gmail SMTP com STARTTLS (porta 587), individualmente para cada
destinatário da lista gerenciada no painel. Ambos os canais são configuráveis sem
reinicialização do contêiner.

\section{Frontend e Painel Web}
\label{sec:frontend}

O \textit{frontend} é uma SPA (\textit{Single Page Application}) desenvolvida em
React 19 com Vite e TailwindCSS para estilização.
A interface exibe dados em sete telas: Dashboard, Detalhes, Métricas, Logs, Alertas,
Endpoints e Configurações. No ambiente implantado, o \textit{frontend} é servido pelo
servidor de desenvolvimento do Vite (porta 5173, interna ao Docker), e não por um
\textit{build} estático otimizado servido por um servidor web, decisão que simplifica
a implantação no contexto acadêmico, mas é registrada como limitação no
Capítulo~\ref{cap:limitacoes}.

O fluxo de \textit{login} envia \texttt{POST /api/auth/login} com a senha do painel;
o servidor devolve a chave de API, que é armazenada no \texttt{localStorage} do
navegador e incluída como cabeçalho \texttt{X-API-Key} em todas as requisições
subsequentes.

As telas de Logs, Alertas e Métricas realizam atualizações automáticas periódicas
por meio de \texttt{setInterval}, a cada 10, 15 e 10 segundos, respectivamente. O
gráfico histórico de métricas e as demais telas são carregados no momento da
navegação ou da troca de \textit{host} selecionado.

As Figuras~\ref{fig:tela-endpoints} e~\ref{fig:tela-config} ilustram duas dessas telas. A tela de Endpoints lista os hosts monitorados com seus identificadores, endereço IP primário e estado de conexão (\textit{online}/\textit{offline}). A tela de Configurações reúne os ajustes do sistema: os limiares de alerta de CPU, memória e disco, os interruptores dos canais de notificação (Teams e e-mail) e a lista de destinatários.

\begin{figure}[htb]
\legenda[fig:tela-endpoints]{Tela de Endpoints com a relação dos hosts monitorados}
\fig{width=1.00\textwidth}{../img/tela-endpoints.png}
\fonte{Os autores (2026).}
\end{figure}

\begin{figure}[htb]
\legenda[fig:tela-config]{Tela de Configurações: limiares de alerta, canais de notificação e destinatários}
\fig{width=1.00\textwidth}{../img/tela-config-1.png}
\fonte{Os autores (2026).}
\end{figure}

\section{Banco de Dados}
\label{sec:banco}

Esta seção descreve a estrutura de persistência do sistema, apresentando o modelo conceitual (entidade-relacionamento), o modelo lógico com as tabelas e seus respectivos tipos de dados, e as estratégias de consulta e de retenção adotadas no PostgreSQL.

\subsection{Modelo Conceitual}
\label{subsec:modelo-conceitual}

O banco de dados é composto por sete tabelas: \texttt{host}, \texttt{agents},
\texttt{host\_discovery}, \texttt{metrics}, \texttt{logs}, \texttt{alerts} e
\texttt{app\_settings}. A tabela \texttt{host} é a entidade central: relaciona-se em
1:1 com \texttt{host\_discovery} e \texttt{agents} e em 1:N com \texttt{metrics},
\texttt{logs} e \texttt{alerts}. A tabela \texttt{app\_settings} é global, armazenando
configurações do sistema (chave de API, destinatários de e-mail, limiares de alerta e
\textit{toggles} de notificação) independentemente de \textit{host}. A
Figura~\ref{fig:modelo-conceitual} apresenta o modelo conceitual do banco no formato
entidade-relacionamento.

\begin{figure}[htb]
\legenda[fig:modelo-conceitual]{Modelo conceitual (entidade-relacionamento) do banco de dados}
\fig{width=1.00\textwidth}{../img/modelo-conceitual-db.png}
\fonte{Os autores (2026).}
\end{figure}

\subsection{Modelo Lógico}
\label{subsec:modelo-logico}

A Figura~\ref{fig:modelo-logico} apresenta o modelo lógico do banco, detalhando as
tabelas, suas colunas e os respectivos tipos de dados.

\begin{figure}[htb]
\legenda[fig:modelo-logico]{Modelo lógico do banco de dados}
\fig{width=1.00\textwidth}{../img/modelo-logico-db.png}
\fonte{Os autores (2026).}
\end{figure}

A tabela \texttt{metrics} registra 20 colunas de séries temporais: além de
\texttt{id}, \texttt{host\_id} e \texttt{timestamp}, armazena percentuais de CPU,
memória e disco; volumes de memória e disco em megabytes; bytes de rede enviados e
recebidos; taxas de leitura e escrita em IOPS e bytes por segundo; e taxas de rede
por segundo. O campo \texttt{parsed\_data} da tabela \texttt{logs} é do tipo JSONB,
permitindo armazenar os campos estruturados extraídos do \texttt{auth.log} sem
colunas fixas adicionais.

\subsection{Consultas e Política de Retenção}
\label{subsec:retencao}

O banco contém 20 índices, incluindo índices GIN sobre os campos JSONB de
\textit{discovery} e logs, e índices parciais sobre \texttt{logs.parsed\_data} para
otimizar as consultas de detecção de força bruta (filtro por
\texttt{status = 'failed'} e \texttt{ip\_origem}).

A retenção de dados é gerenciada pela função PL/pgSQL
\texttt{cleanup\_old\_data(days\_to\_keep INT DEFAULT 7)}, definida no esquema
inicial, e executada automaticamente pelo APScheduler diariamente às 03:00 UTC.
A chamada ocorre via \texttt{SELECT cleanup\_old\_data(:dias)} no banco. A função
elimina registros das tabelas \texttt{metrics} e \texttt{logs} com
\texttt{timestamp} anterior ao período de retenção configurado. A limpeza também
pode ser acionada manualmente via \texttt{POST /api/maintenance/cleanup}.

% =================================================================================================
% 13. RESULTADOS E VALIDAÇÃO
% =================================================================================================
\chapter{Resultados e Validação}
\label{cap:resultados}

Este capítulo apresenta a validação funcional do sistema implementado, conduzida em ambiente de rede local controlado. Cada seção verifica um componente específico da solução: a coleta de métricas, o inventário de \textit{hardware} (\textit{discovery}), a coleta de \textit{logs} de autenticação, a detecção de força bruta e de varredura de portas, e a geração de alertas com as respectivas notificações. O objetivo é demonstrar, por meio de evidências capturadas durante a execução, que o sistema opera de ponta a ponta conforme especificado.

\section{Ambiente de Teste}
\label{sec:ambiente-teste}

A validação foi conduzida no mesmo ambiente de rede local descrito no Capítulo~\ref{cap:implementacao}: um servidor central e seis clientes Linux conectados em topologia estrela. Para os testes de segurança, utilizou-se um host atacante na mesma sub-rede, executando as ferramentas \texttt{hydra} (força bruta SSH) e \texttt{nmap} (varredura de portas). O objetivo desta etapa foi verificar, de forma funcional, que cada componente opera conforme especificado e que a solução funciona de ponta a ponta, e não a aferição de métricas de desempenho. A Figura~\ref{fig:val-dashboard} apresenta a visão geral do painel, com os hosts monitorados. Cabe um esclarecimento sobre os endereços exibidos: cada máquina de teste possuía duas interfaces de rede: uma conectada ao \textit{switch} da rede de monitoramento (\texttt{10.10.10.0/26}) e outra à rede doméstica com acesso à internet (\texttt{192.168.0.0/24}), utilizada para baixar o agente. Como o agente reporta o IP primário do host, isto é, o associado à rota padrão (\textit{default route}) obtida via \texttt{ip route}, o endereço apresentado no painel corresponde à interface de internet (\texttt{192.168.0.x}). O tráfego de ataque (força bruta e varredura), por sua vez, percorre a rede de monitoramento, o que explica o IP atacante \texttt{10.10.10.x} observado nos alertas.

\begin{figure}[htb]
\legenda[fig:val-dashboard]{Painel principal exibindo os hosts monitorados}
\fig{width=1.00\textwidth}{../img/dashboard-pos-instalacao.png}
\fonte{Os autores (2026).}
\end{figure}

\section{Coleta de Métricas}
\label{sec:valid-metricas}

Com o agente instalado e ativo em um cliente, verificou-se no painel a exibição contínua das métricas de CPU, memória e disco, atualizadas automaticamente. Os valores reportados pelo agente foram conferidos contra as ferramentas nativas do sistema operacional (\texttt{top}, \texttt{free -m} e \texttt{df -h}), confirmando que correspondem ao estado real do host monitorado. As Figuras~\ref{fig:val-metricas-1} e~\ref{fig:val-metricas-2} apresentam a tela de métricas: a primeira exibe os medidores de status em tempo real de CPU, memória e disco, acompanhados do gráfico de evolução; a segunda detalha o histórico das métricas e os cartões com o uso atual de cada recurso.

\begin{figure}[htb]
\legenda[fig:val-metricas-1]{Tela de métricas: medidores em tempo real e gráfico de evolução}
\fig{width=1.00\textwidth}{../img/metricas-host-1.png}
\fonte{Os autores (2026).}
\end{figure}

\begin{figure}[htb]
\legenda[fig:val-metricas-2]{Tela de métricas: histórico detalhado e cartões de uso atual}
\fig{width=1.00\textwidth}{../img/metricas-host-2.png}
\fonte{Os autores (2026).}
\end{figure}

\section{Coleta de Inventário (Discovery): Físico \textit{versus} Virtual}
\label{sec:valid-discovery}

Na primeira execução, o agente realiza a coleta de inventário (\textit{discovery}) do host, distinguindo máquinas físicas de virtuais conforme a lógica descrita na Seção~\ref{subsec:discovery}. As Figuras~\ref{fig:val-vhost-1} e~\ref{fig:val-vhost-2} apresentam os detalhes de um host virtual, e as Figuras~\ref{fig:val-fisico-1} a~\ref{fig:val-fisico-3}, os de um host físico. Observa-se que o host físico expõe informações de hardware mais detalhadas (modelo de CPU, \textit{clock}, discos, slots e fabricante dos módulos de memória e dados da placa-mãe), enquanto no host virtual parte desses dados é omitida ou reportada de forma genérica, comportamento esperado, já que o ambiente virtualizado não dá acesso ao hardware real.

\begin{figure}[htb]
\legenda[fig:val-vhost-1]{Host virtual: visão geral de CPU, memória, disco e sistema}
\fig{width=1.00\textwidth}{../img/detalhes-vhost-1.png}
\fonte{Os autores (2026).}
\end{figure}

\begin{figure}[htb]
\legenda[fig:val-vhost-2]{Host virtual: discos e interfaces de rede}
\fig{width=1.00\textwidth}{../img/detalhes-vhost-2.png}
\fonte{Os autores (2026).}
\end{figure}

\begin{figure}[htb]
\legenda[fig:val-fisico-1]{Host físico: visão geral de hardware e sistema}
\fig{width=1.00\textwidth}{../img/detalhes-hostFisico-1.png}
\fonte{Os autores (2026).}
\end{figure}

\begin{figure}[htb]
\legenda[fig:val-fisico-2]{Host físico: detalhes de discos e interfaces de rede}
\fig{width=1.00\textwidth}{../img/detalhes-hostFisico-2.png}
\fonte{Os autores (2026).}
\end{figure}

\begin{figure}[htb]
\legenda[fig:val-fisico-3]{Host físico: módulos de memória RAM e placa-mãe}
\fig{width=1.00\textwidth}{../img/detalhes-hostFisico-3.png}
\fonte{Os autores (2026).}
\end{figure}

\section{Coleta de Logs de Autenticação}
\label{sec:valid-logs}

Para validar a coleta de logs, foram realizadas tentativas de acesso SSH ao cliente monitorado, uma malsucedida e uma bem-sucedida. Ambas foram lidas pelo agente a partir do arquivo \texttt{/var/log/auth.log} e exibidas na tela de Logs do painel, com os campos relevantes (origem, usuário e resultado) corretamente extraídos. A Figura~\ref{fig:val-logs} mostra os eventos registrados.

\begin{figure}[htb]
\legenda[fig:val-logs]{Tela de Logs exibindo os eventos de autenticação coletados}
\fig{width=1.00\textwidth}{../img/logs-gerais.png}
\fonte{Os autores (2026).}
\end{figure}

\section{Detecção de Força Bruta SSH}
\label{sec:valid-bruteforce}

Simulou-se um ataque de força bruta com a ferramenta \texttt{hydra}, gerando múltiplas tentativas de autenticação SSH malsucedidas a partir do host atacante. Ao ultrapassar o limiar de cinco falhas do mesmo endereço em uma janela de 60 segundos, o \textit{backend} gerou automaticamente um alerta do tipo \texttt{brute\_force} e o encaminhou aos canais de notificação. A Figura~\ref{fig:val-bf-cmd} mostra o comando \texttt{hydra} em execução contra o cliente; as Figuras~\ref{fig:val-bf-teams} e~\ref{fig:val-bf-email} apresentam, respectivamente, a notificação do alerta recebida no Microsoft \textit{Teams} e por e-mail.

\begin{figure}[htb]
\legenda[fig:val-bf-cmd]{Execução do ataque de força bruta com \texttt{hydra}}
\fig{width=1.00\textwidth}{../img/brute-force-1-comandoHYDRA.png}
\fonte{Os autores (2026).}
\end{figure}

\begin{figure}[htb]
\legenda[fig:val-bf-teams]{Notificação do alerta de força bruta no Microsoft Teams}
\fig{width=0.65\textwidth}{../img/brute-force-2-alertaTeams.png}
\fonte{Os autores (2026).}
\end{figure}

\begin{figure}[htb]
\legenda[fig:val-bf-email]{Notificação do alerta de força bruta por e-mail}
\fig{width=0.85\textwidth}{../img/brute-force-3-alertaEmail.png}
\fonte{Os autores (2026).}
\end{figure}

\section{Detecção de Port Scan}
\label{sec:valid-portscan}

A detecção de varredura de portas foi validada executando-se o \texttt{nmap} contra o cliente monitorado. O agente, por meio da captura de pacotes com \texttt{tcpdump}, identificou o acesso a múltiplas portas distintas a partir de um mesmo endereço e sinalizou o evento ao \textit{backend}, que gerou um alerta do tipo \texttt{port\_scan}. A Figura~\ref{fig:val-ps-cmd} mostra o \texttt{nmap} em execução; a Figura~\ref{fig:val-ps-painel} apresenta o alerta de \texttt{port\_scan} no painel; e as Figuras~\ref{fig:val-ps-email} e~\ref{fig:val-ps-teams} mostram as notificações correspondentes por e-mail e no \textit{Teams}. Vale observar que o \texttt{nmap}, em sua varredura padrão, acessou cerca de mil portas distintas (número muito acima do limiar de detecção de dez portas distintas em 60 segundos definido na Seção~\ref{subsec:portscan-agente}), de modo que o disparo do alerta era esperado.

\begin{figure}[htb]
\legenda[fig:val-ps-cmd]{Execução da varredura de portas com \texttt{nmap}}
\fig{width=1.00\textwidth}{../img/port-scan-1-comandoNMAP.png}
\fonte{Os autores (2026).}
\end{figure}

\begin{figure}[htb]
\legenda[fig:val-ps-painel]{Alerta de varredura de portas exibido no painel}
\fig{width=1.00\textwidth}{../img/port-scan-2-alertaGeral.png}
\fonte{Os autores (2026).}
\end{figure}

\begin{figure}[htb]
\legenda[fig:val-ps-email]{Notificação do alerta de varredura de portas por e-mail}
\fig{width=0.85\textwidth}{../img/port-scan-3-alertaEmail.png}
\fonte{Os autores (2026).}
\end{figure}

\begin{figure}[htb]
\legenda[fig:val-ps-teams]{Notificação do alerta de varredura de portas no Microsoft Teams}
\fig{width=0.65\textwidth}{../img/port-scan-4-alertaTeams.png}
\fonte{Os autores (2026).}
\end{figure}

\section{Geração de Alertas e Notificações}
\label{sec:valid-notificacoes}

Além dos alertas de segurança apresentados nas seções anteriores, o sistema gera alertas de recurso quando o uso de CPU, memória ou disco de um host ultrapassa o limiar configurado. Para validar esse mecanismo, induziu-se carga em um cliente até que o uso de CPU e de memória excedesse os limiares definidos, observando-se a geração automática dos alertas correspondentes.

Todos os alertas, de segurança e de recurso, são consolidados na tela de Alertas Gerais do painel. As Figuras~\ref{fig:val-alertas-1} e~\ref{fig:val-alertas-2} apresentam essa tela: primeiro com um único alerta de CPU e, em seguida, com múltiplos alertas de tipos distintos (força bruta, CPU e varredura de portas) registrados simultaneamente.

Com os canais habilitados, cada alerta é encaminhado por e-mail e ao Microsoft \textit{Teams}. As Figuras~\ref{fig:val-cpu-email} e~\ref{fig:val-cpu-teams} mostram a notificação de um alerta de CPU alta nos dois canais, e as Figuras~\ref{fig:val-mem-email} e~\ref{fig:val-mem-teams}, a de um alerta de memória alta. Verificou-se também que, ao desabilitar um canal nas configurações, a respectiva notificação deixa de ser enviada, embora o alerta continue registrado no painel.

\begin{figure}[htb]
\legenda[fig:val-alertas-1]{Tela de Alertas Gerais com um alerta de recurso (CPU)}
\fig{width=1.00\textwidth}{../img/alertas-gerais.png}
\fonte{Os autores (2026).}
\end{figure}

\begin{figure}[htb]
\legenda[fig:val-alertas-2]{Tela de Alertas Gerais com múltiplos alertas de tipos distintos}
\fig{width=1.00\textwidth}{../img/alertas-gerais-2-comMaisAlertas.png}
\fonte{Os autores (2026).}
\end{figure}

\begin{figure}[htb]
\legenda[fig:val-cpu-email]{Notificação de alerta de CPU alta por e-mail}
\fig{width=0.85\textwidth}{../img/alertaEmail-CPU.png}
\fonte{Os autores (2026).}
\end{figure}

\begin{figure}[htb]
\legenda[fig:val-cpu-teams]{Notificação de alerta de CPU alta no Microsoft Teams}
\fig{width=0.65\textwidth}{../img/alertaTeams-CPU.png}
\fonte{Os autores (2026).}
\end{figure}

\begin{figure}[htb]
\legenda[fig:val-mem-email]{Notificação de alerta de memória alta por e-mail}
\fig{width=0.85\textwidth}{../img/alertaEmail-Memoria.png}
\fonte{Os autores (2026).}
\end{figure}

\begin{figure}[htb]
\legenda[fig:val-mem-teams]{Notificação de alerta de memória alta no Microsoft Teams}
\fig{width=0.65\textwidth}{../img/alertaTeams-Memoria.png}
\fonte{Os autores (2026).}
\end{figure}

\section{Síntese da Validação}
\label{sec:discussao}

Os testes confirmam que o sistema atende aos objetivos definidos no Capítulo~\ref{cap:introducao}: o agente coleta métricas de infraestrutura, inventário de \textit{hardware} e logs de autenticação; o \textit{backend} processa esses dados, detecta tentativas de força bruta e varreduras de portas e gera os alertas correspondentes; e o painel apresenta as informações ao operador, que ainda é notificado por \textit{Teams} e e-mail. A validação teve caráter funcional, demonstrando o funcionamento de ponta a ponta da solução em ambiente real; a aferição quantitativa de desempenho, como latência de detecção e sobrecarga do agente, não fez parte do escopo deste trabalho.

% =================================================================================================
% 14. LIMITAÇÕES E TRABALHOS FUTUROS
% =================================================================================================
\chapter{Limitações e Trabalhos Futuros}
\label{cap:limitacoes}

Este capítulo discute as limitações da solução desenvolvida, decorrentes das escolhas arquiteturais e do escopo acadêmico do projeto, e apresenta, a partir delas, um conjunto de trabalhos futuros que apontam caminhos para a evolução do sistema.

\section{Limitações da Solução}
\label{sec:limitacoes}

A implementação do sistema revelou um conjunto de limitações que decorrem de escolhas
arquiteturais e do escopo do projeto:

\begin{itemize}
    \item \textbf{Ausência de comunicação segura (TLS/HTTPS)}: toda a comunicação
          entre agentes, \textit{frontend} e \textit{backend} ocorre em HTTP puro,
          sem cifragem. A arquitetura original previa o uso de Caddy como
          \textit{proxy} reverso com Step-CA para emissão de certificados internos;
          essa funcionalidade foi deliberadamente adiada para reduzir a complexidade
          de implantação no ambiente acadêmico.
    \item \textbf{Ponto único de falha}: a arquitetura centralizada, com um único
          servidor de \textit{backend} e banco de dados sem réplica, implica que uma
          falha no servidor interrompe o monitoramento de todos os clientes.
    \item \textbf{Ausência de controle de acesso granular}: uma única chave de API
          concede acesso irrestrito a todos os \textit{endpoints} protegidos. Não
          há suporte a múltiplos usuários, perfis ou permissões diferenciadas.
    \item \textbf{Sem limitação de taxa (\textit{rate limiting})}: o \textit{backend}
          não implementa restrições de frequência de requisições.
    \item \textbf{Port scan sem persistência}: os eventos de varredura detectados pelo
          agente não são persistidos em tabela própria; o sistema registra apenas o
          alerta gerado, sem histórico de portas escaneadas.
    \item \textbf{Limiar de força bruta não configurável}: o parâmetro
          \texttt{LIMIAR\_BRUTE\_FORCE} (5 tentativas) está definido como constante
          no código-fonte, diferentemente dos limiares de CPU, memória e disco que são
          configuráveis via painel.
    \item \textbf{\textit{Frontend} em modo de desenvolvimento}: o contêiner do
          \textit{frontend} executa o servidor de desenvolvimento do Vite
          (\texttt{NODE\_ENV=development}, com observadores de arquivos ativos), em vez
          de um \textit{build} estático otimizado servido por um servidor web. A escolha
          simplifica a implantação, mas não é adequada a um ambiente de produção.
    \item \textbf{Dependência do arquivo \texttt{/var/log/auth.log}}: a coleta de logs
          de autenticação e, por consequência, a detecção de força bruta dependem da
          existência do arquivo \texttt{/var/log/auth.log}, populado por um \textit{daemon}
          de \textit{syslog} (como o \texttt{rsyslog}). Em distribuições configuradas apenas
          com o \texttt{systemd-journald}, sem \texttt{rsyslog} (situação observada em um
          cliente Kali Linux durante os testes), esse arquivo não é gerado e tais
          funcionalidades ficam indisponíveis no host. As demais funcionalidades (métricas,
          \textit{discovery} e detecção de \textit{port scan}, que não dependem do
          \texttt{auth.log}) permanecem operacionais.
\end{itemize}

\section{Trabalhos Futuros}
\label{sec:futuros}

Com base nas limitações identificadas e nos objetivos de extensão da plataforma, os
seguintes trabalhos futuros são propostos:

\begin{itemize}
    \item \textbf{Adoção de TLS/HTTPS}: implementar a arquitetura originalmente
          planejada com Caddy como \textit{proxy} reverso e Step-CA como autoridade
          certificadora interna, cifrando toda a comunicação.
    \item \textbf{Persistência dos eventos de port scan}: criar tabela dedicada para
          registrar o histórico de varreduras por IP e porta, permitindo análise
          temporal e correlação com outros eventos.
    \item \textbf{Controle de acesso granular (RBAC)}: implementar autenticação
          baseada em papéis, com usuários e permissões diferenciadas.
    \item \textbf{Configuração dinâmica do limiar de força bruta}: mover o parâmetro
          para a tabela \texttt{app\_settings}, tornando-o configurável pelo painel.
    \item \textbf{Detecção por anomalia}: complementar as regras estáticas por
          modelos de linha de base que detectem desvios comportamentais não
          contemplados por limiares fixos, conforme discutido em
          \citeonline{sbseg2023}.
    \item \textbf{Arquitetura orientada a eventos}: introduzir um
          \textit{message broker} para desacoplar a ingestão de dados do
          processamento, aumentando resiliência e escalabilidade conforme
          \citeonline{aws2025eda}.
    \item \textbf{Alta disponibilidade}: adicionar réplica de leitura ao PostgreSQL
          e balanceamento de carga ao \textit{backend}, eliminando o ponto único de
          falha.
    \item \textbf{Correção das lacunas do frontend}: resolver os comportamentos
          pendentes identificados no painel web.
    \item \textbf{Leitura do \textit{journal} do \textit{systemd} e detecção de
          distribuição}: ler o \texttt{systemd-journald} (via \texttt{journalctl}) quando o
          \texttt{/var/log/auth.log} não existir e identificar a distribuição do host para
          adaptar as fontes de \textit{log}, removendo a dependência de um \textit{daemon}
          de \textit{syslog}.
\end{itemize}

% =================================================================================================
% 15. USO DE INTELIGÊNCIA ARTIFICIAL NO TRABALHO
% NOTA: confirmar com orientador se vai como capítulo numerado ou apêndice.
% =================================================================================================
\chapter{Uso de Inteligência Artificial no Trabalho}
\label{cap:ia}

Durante o desenvolvimento deste projeto, a equipe fez uso de uma ferramenta de
inteligência artificial generativa como recurso de apoio. Este capítulo descreve
de forma transparente o escopo e os limites desse uso, em conformidade com as
diretrizes de integridade acadêmica.

\section{Ferramenta Utilizada}
\label{sec:ia-ferramenta}

Foi utilizado o assistente Claude (Anthropic), acessado entre maio e junho de 2026
por meio da interface web e da interface de linha de comando Claude Code. Nenhuma
outra ferramenta de inteligência artificial generativa foi empregada no projeto.

\section{Finalidades do Uso}
\label{sec:ia-finalidades}

O assistente foi utilizado exclusivamente nas seguintes finalidades:

\begin{itemize}
    \item \textbf{Compreensão de código}: análise e explicação de trechos do
          código-fonte implementado pela equipe, auxiliando no entendimento de
          comportamentos e interações entre componentes;
    \item \textbf{Resolução de problemas técnicos}: apoio na identificação de causas
          de erros e na busca de abordagens de solução, sempre verificadas pela
          equipe contra o código real;
    \item \textbf{Validação de ideias}: discussão de decisões de projeto para
          verificar consistência lógica e identificar possíveis lacunas, antes de
          a equipe tomar a decisão final;
    \item \textbf{Sugestão de caminhos de investigação}: indicação de arquivos e
          funções relevantes a inspecionar durante o processo de desenvolvimento e
          documentação;
    \item \textbf{Redação de base e revisão do texto}: geração de versões iniciais
          (rascunhos) de trechos deste documento e apoio na revisão de redação,
          consistência e formatação, posteriormente reescritos e ajustados pela equipe.
\end{itemize}

\section{Responsabilidade e Autoria}
\label{sec:ia-responsabilidade}

Todas as decisões de arquitetura, infraestrutura e implementação foram tomadas
integralmente pela equipe. Nenhuma decisão técnica foi delegada à ferramenta de
inteligência artificial. Os trechos de texto elaborados com o auxílio do assistente
serviram como base inicial e foram posteriormente reescritos e ajustados pelos autores,
de forma que a autoria final do documento é da equipe. Todo o conteúdo foi revisado
criticamente e verificado contra o código-fonte do repositório antes de ser incorporado
a este documento. Os autores assumem responsabilidade integral pelo conteúdo aqui
apresentado.


% =================================================================================================
% 16. CONSIDERAÇÕES FINAIS
% =================================================================================================
\chapter{Considerações Finais}
\label{cap:conclusao}

Este trabalho teve como objetivo desenvolver um sistema de monitoramento para
ambientes Linux, capaz de coletar métricas operacionais e identificar eventos de
segurança. Ao longo dos capítulos anteriores, apresentou-se a fundamentação teórica
que embasou as decisões de projeto, o planejamento arquitetural e a implementação
concreta da solução.

O objetivo geral foi plenamente atendido. O sistema foi implementado, integrado e
operado em rede local com múltiplos clientes Linux reais. Os quatro objetivos
específicos definidos no Capítulo~\ref{cap:introducao} foram alcançados: (i)
mecanismos de coleta foram implementados no agente, com ciclos de 5 segundos e
\textit{discovery} de \textit{hardware}; (ii) o servidor central processa e armazena
os dados via API REST e PostgreSQL; (iii) a detecção de força bruta SSH e varredura
de portas foi implementada; e (iv) o painel \textit{web} oferece visualização em
tempo quase real de métricas, \textit{logs} e alertas.

A evolução do sistema ao longo do desenvolvimento evidenciou a importância do processo
iterativo. Três decisões merecem destaque: a migração da detecção de
\textit{port scan} do \textit{backend} para o agente (com \texttt{tcpdump}), que
passou a operar diretamente sobre a camada de rede, capturando pacotes SYN, em vez de
inferir varredura pela frequência de conexões registradas; a remoção da tabela \texttt{active\_connections} após constatar
que seu custo operacional superava o valor analítico; e o endurecimento da
autenticação, que substituiu um \textit{token} estático por uma chave de API dinâmica
de 256 bits.

A fundamentação teórica mostrou-se diretamente aplicável. O método USE orientou a
seleção das métricas coletadas. Os conceitos de observabilidade reforçaram a
importância de combinar métricas, \textit{logs} e alertas em uma visão integrada.
As diretrizes de \citeonline{nist_sp800_92} fundamentaram a coleta e análise de
\textit{logs} de autenticação. A arquitetura cliente-servidor com agentes autônomos,
conforme descrita por \citeonline{wooldridge1995}, provou-se adequada ao contexto de
monitoramento distribuído em rede local.

O sistema apresenta limitações reais, discutidas no Capítulo~\ref{cap:limitacoes},
com destaque para a ausência de comunicação segura (HTTP puro) e o ponto único de
falha na arquitetura centralizada. Essas limitações refletem o escopo acadêmico do
projeto e apontam caminhos claros para evolução futura.

Dessa forma, o trabalho cumpre seu papel ao integrar pesquisa teórica, planejamento e
implementação prática de um sistema de monitoramento funcional, contribuindo como
referência para projetos similares em ambientes acadêmicos e de pequeno porte.

\postextual

\bibliography{./referencias}

\end{document}